% Môn: Vật lý
% Lớp: 12
% Chương: Chương IV. Vật lí hạt nhân
% Bài: L12C4 Bài 23. Hiện tượng phóng xạ
% Số câu: 162
% App đọc file này trực tiếp — sửa rồi Commit trên GitHub, bấm Đồng bộ GitHub .tex

% L12C4 Bài 23. Hiện tượng phóng xạ

% ===== Câu 1 =====
% ID: L12C4B23-01-DS
% Mức: TH
\dangbt{Nhận biết hiện tượng phóng xạ và các tia phóng xạ}
\begin{ex}
% ID: L12C4B23-01-DS
% Mức: TH
Xét nhận biết hiện tượng phóng xạ và các tia phóng xạ (tình huống 1). Hãy xác định tính đúng, sai của bốn phát biểu sau.
\choiceTF
{\True Phóng xạ là quá trình tự phát của hạt nhân không bền; tia $\alpha,\beta,\gamma$ có bản chất và khả năng đâm xuyên khác nhau.}
{Cường độ phóng xạ có thể tăng tùy ý bằng cách nung nóng mẫu.}
{\True Khi giải cần đọc đúng dữ kiện, dùng hệ đơn vị thống nhất và kiểm tra điều kiện áp dụng.}
{Có thể bỏ qua dữ kiện, đơn vị và ý nghĩa vật lí của kết quả.}
\loigiai{
\begin{itemchoice}
\item (Đúng) Phóng xạ là quá trình tự phát của hạt nhân không bền; tia $\alpha,\beta,\gamma$ có bản chất và khả năng đâm xuyên khác nhau. (Vì): Phóng xạ là quá trình tự phát của hạt nhân không bền; tia $\alpha,\beta,\gamma$ có bản chất và khả năng đâm xuyên khác nhau. b) Sai: Cường độ phóng xạ có thể tăng tùy ý bằng cách nung nóng mẫu. c) Đúng: Khi giải cần đọc đúng dữ kiện, dùng hệ đơn vị thống nhất và kiểm tra điều kiện áp dụng. d) Sai: Có thể bỏ qua dữ kiện, đơn vị và ý nghĩa vật lí của kết quả. Kiến thức cốt lõi: Phóng xạ là quá trình tự phát của hạt nhân không bền; tia $\alpha,\beta,\gamma$ có bản chất và khả năng đâm xuyên khác nhau
\item (Sai) Cường độ phóng xạ có thể tăng tùy ý bằng cách nung nóng mẫu.
\item (Đúng) Khi giải cần đọc đúng dữ kiện, dùng hệ đơn vị thống nhất và kiểm tra điều kiện áp dụng.
\item (Sai) Có thể bỏ qua dữ kiện, đơn vị và ý nghĩa vật lí của kết quả.
\end{itemchoice}
}
\end{ex}

% ===== Câu 2 =====
% ID: L12C4B23-01-TL
% Mức: TH
\begin{bt}
% ID: L12C4B23-01-TL
% Mức: TH
Trình bày kiến thức cốt lõi của nhận biết hiện tượng phóng xạ và các tia phóng xạ và nêu điều kiện áp dụng.
\loigiai{
Cần nêu được: Phóng xạ là quá trình tự phát của hạt nhân không bền; tia $\alpha,\beta,\gamma$ có bản chất và khả năng đâm xuyên khác nhau. Khi vận dụng phải xác định đúng dữ kiện, điều kiện áp dụng, hệ đơn vị và kiểm tra tính hợp lí của kết quả. Nhận định “Cường độ phóng xạ có thể tăng tùy ý bằng cách nung nóng mẫu.” sai vì trái với kiến thức trên.
}
\end{bt}

% ===== Câu 3 =====
% ID: L12C4B23-01-TLN
% Mức: TH
\begin{ex}
% ID: L12C4B23-01-TLN
% Mức: TH
Trong bốn phát biểu sau về nhận biết hiện tượng phóng xạ và các tia phóng xạ, có bao nhiêu phát biểu đúng? Chỉ ghi một số nguyên. (1) Phóng xạ là quá trình tự phát của hạt nhân không bền; tia $\alpha,\beta,\gamma$ có bản chất và khả năng đâm xuyên khác nhau. (2) Cường độ phóng xạ có thể tăng tùy ý bằng cách nung nóng mẫu. (3) Cần kiểm tra điều kiện áp dụng trước khi tính toán. (4) Có thể bỏ qua đơn vị của kết quả.
\shortans{2}
\loigiai{
Đối chiếu từng phát biểu với kiến thức cốt lõi: Phóng xạ là quá trình tự phát của hạt nhân không bền; tia $\alpha,\beta,\gamma$ có bản chất và khả năng đâm xuyên khác nhau. Có 2 phát biểu đúng.
}
\end{ex}

% ===== Câu 4 =====
% ID: L12C4B23-01-TN
% Mức: NB
\begin{ex}
% ID: L12C4B23-01-TN
% Mức: NB
Phát biểu nào sau đây đúng về nhận biết hiện tượng phóng xạ và các tia phóng xạ?
\choice
{\True Phóng xạ là quá trình tự phát của hạt nhân không bền; tia $\alpha,\beta,\gamma$ có bản chất và khả năng đâm xuyên khác nhau.}
{Cường độ phóng xạ có thể tăng tùy ý bằng cách nung nóng mẫu.}
{Có thể bỏ qua điều kiện áp dụng và đơn vị trong mọi trường hợp.}
{Kết quả của bài toán không phụ thuộc dữ kiện đã cho.}
\loigiai{
Phóng xạ là quá trình tự phát của hạt nhân không bền; tia $\alpha,\beta,\gamma$ có bản chất và khả năng đâm xuyên khác nhau. Vì vậy phương án $A$ đúng; các phương án còn lại trái với kiến thức hoặc điều kiện áp dụng.
}
\end{ex}

% ===== Câu 5 =====
% ID: L12C4B23-02-DS
% Mức: TH
\begin{ex}
% ID: L12C4B23-02-DS
% Mức: TH
Xét nhận biết hiện tượng phóng xạ và các tia phóng xạ (tình huống 2). Hãy xác định tính đúng, sai của bốn phát biểu sau.
\choiceTF
{Cường độ phóng xạ có thể tăng tùy ý bằng cách nung nóng mẫu.}
{\True Phóng xạ là quá trình tự phát của hạt nhân không bền; tia $\alpha,\beta,\gamma$ có bản chất và khả năng đâm xuyên khác nhau.}
{Có thể bỏ qua dữ kiện, đơn vị và ý nghĩa vật lí của kết quả.}
{\True Khi giải cần đọc đúng dữ kiện, dùng hệ đơn vị thống nhất và kiểm tra điều kiện áp dụng.}
\loigiai{
\begin{itemchoice}
\item (Sai) Cường độ phóng xạ có thể tăng tùy ý bằng cách nung nóng mẫu. (Vì): Cường độ phóng xạ có thể tăng tùy ý bằng cách nung nóng mẫu. b) Đúng: Phóng xạ là quá trình tự phát của hạt nhân không bền; tia $\alpha,\beta,\gamma$ có bản chất và khả năng đâm xuyên khác nhau. c) Sai: Có thể bỏ qua dữ kiện, đơn vị và ý nghĩa vật lí của kết quả. d) Đúng: Khi giải cần đọc đúng dữ kiện, dùng hệ đơn vị thống nhất và kiểm tra điều kiện áp dụng. Kiến thức cốt lõi: Phóng xạ là quá trình tự phát của hạt nhân không bền; tia $\alpha,\beta,\gamma$ có bản chất và khả năng đâm xuyên khác nhau
\item (Đúng) Phóng xạ là quá trình tự phát của hạt nhân không bền; tia $\alpha,\beta,\gamma$ có bản chất và khả năng đâm xuyên khác nhau.
\item (Sai) Có thể bỏ qua dữ kiện, đơn vị và ý nghĩa vật lí của kết quả.
\item (Đúng) Khi giải cần đọc đúng dữ kiện, dùng hệ đơn vị thống nhất và kiểm tra điều kiện áp dụng.
\end{itemchoice}
}
\end{ex}

% ===== Câu 6 =====
% ID: L12C4B23-02-TL
% Mức: TH
\begin{bt}
% ID: L12C4B23-02-TL
% Mức: TH
Giải thích vì sao nhận định sau đúng: “Phóng xạ là quá trình tự phát của hạt nhân không bền; tia $\alpha,\beta,\gamma$ có bản chất và khả năng đâm xuyên khác nhau.”
\loigiai{
Cần nêu được: Phóng xạ là quá trình tự phát của hạt nhân không bền; tia $\alpha,\beta,\gamma$ có bản chất và khả năng đâm xuyên khác nhau. Khi vận dụng phải xác định đúng dữ kiện, điều kiện áp dụng, hệ đơn vị và kiểm tra tính hợp lí của kết quả. Nhận định “Cường độ phóng xạ có thể tăng tùy ý bằng cách nung nóng mẫu.” sai vì trái với kiến thức trên.
}
\end{bt}

% ===== Câu 7 =====
% ID: L12C4B23-02-TLN
% Mức: TH
\begin{ex}
% ID: L12C4B23-02-TLN
% Mức: TH
Trong bốn phát biểu sau về nhận biết hiện tượng phóng xạ và các tia phóng xạ, có bao nhiêu phát biểu đúng? Chỉ ghi một số nguyên. (1) Cường độ phóng xạ có thể tăng tùy ý bằng cách nung nóng mẫu. (2) Phóng xạ là quá trình tự phát của hạt nhân không bền; tia $\alpha,\beta,\gamma$ có bản chất và khả năng đâm xuyên khác nhau. (3) Kết quả phải phù hợp ý nghĩa vật lí. (4) Mọi đại lượng đều có cùng bản chất.
\shortans{2}
\loigiai{
Đối chiếu từng phát biểu với kiến thức cốt lõi: Phóng xạ là quá trình tự phát của hạt nhân không bền; tia $\alpha,\beta,\gamma$ có bản chất và khả năng đâm xuyên khác nhau. Có 2 phát biểu đúng.
}
\end{ex}

% ===== Câu 8 =====
% ID: L12C4B23-02-TN
% Mức: NB
\begin{ex}
% ID: L12C4B23-02-TN
% Mức: NB
Nhận định nào phù hợp với kiến thức Vật lí về nhận biết hiện tượng phóng xạ và các tia phóng xạ?
\choice
{Cường độ phóng xạ có thể tăng tùy ý bằng cách nung nóng mẫu.}
{\True Phóng xạ là quá trình tự phát của hạt nhân không bền; tia $\alpha,\beta,\gamma$ có bản chất và khả năng đâm xuyên khác nhau.}
{Có thể bỏ qua điều kiện áp dụng và đơn vị trong mọi trường hợp.}
{Kết quả của bài toán không phụ thuộc dữ kiện đã cho.}
\loigiai{
Phóng xạ là quá trình tự phát của hạt nhân không bền; tia $\alpha,\beta,\gamma$ có bản chất và khả năng đâm xuyên khác nhau. Vì vậy phương án $B$ đúng; các phương án còn lại trái với kiến thức hoặc điều kiện áp dụng.
}
\end{ex}

% ===== Câu 9 =====
% ID: L12C4B23-03-DS
% Mức: VD
\begin{ex}
% ID: L12C4B23-03-DS
% Mức: VD
Xét nhận biết hiện tượng phóng xạ và các tia phóng xạ (tình huống 3). Hãy xác định tính đúng, sai của bốn phát biểu sau.
\choiceTF
{\True Phóng xạ là quá trình tự phát của hạt nhân không bền; tia $\alpha,\beta,\gamma$ có bản chất và khả năng đâm xuyên khác nhau.}
{Cường độ phóng xạ có thể tăng tùy ý bằng cách nung nóng mẫu.}
{\True Khi giải cần đọc đúng dữ kiện, dùng hệ đơn vị thống nhất và kiểm tra điều kiện áp dụng.}
{Có thể bỏ qua dữ kiện, đơn vị và ý nghĩa vật lí của kết quả.}
\loigiai{
\begin{itemchoice}
\item (Đúng) Phóng xạ là quá trình tự phát của hạt nhân không bền; tia $\alpha,\beta,\gamma$ có bản chất và khả năng đâm xuyên khác nhau. (Vì): Phóng xạ là quá trình tự phát của hạt nhân không bền; tia $\alpha,\beta,\gamma$ có bản chất và khả năng đâm xuyên khác nhau. b) Sai: Cường độ phóng xạ có thể tăng tùy ý bằng cách nung nóng mẫu. c) Đúng: Khi giải cần đọc đúng dữ kiện, dùng hệ đơn vị thống nhất và kiểm tra điều kiện áp dụng. d) Sai: Có thể bỏ qua dữ kiện, đơn vị và ý nghĩa vật lí của kết quả. Kiến thức cốt lõi: Phóng xạ là quá trình tự phát của hạt nhân không bền; tia $\alpha,\beta,\gamma$ có bản chất và khả năng đâm xuyên khác nhau
\item (Sai) Cường độ phóng xạ có thể tăng tùy ý bằng cách nung nóng mẫu.
\item (Đúng) Khi giải cần đọc đúng dữ kiện, dùng hệ đơn vị thống nhất và kiểm tra điều kiện áp dụng.
\item (Sai) Có thể bỏ qua dữ kiện, đơn vị và ý nghĩa vật lí của kết quả.
\end{itemchoice}
}
\end{ex}

% ===== Câu 10 =====
% ID: L12C4B23-03-TL
% Mức: VD
\begin{bt}
% ID: L12C4B23-03-TL
% Mức: VD
Phân tích sai lầm trong nhận định: “Cường độ phóng xạ có thể tăng tùy ý bằng cách nung nóng mẫu.”
\loigiai{
Cần nêu được: Phóng xạ là quá trình tự phát của hạt nhân không bền; tia $\alpha,\beta,\gamma$ có bản chất và khả năng đâm xuyên khác nhau. Khi vận dụng phải xác định đúng dữ kiện, điều kiện áp dụng, hệ đơn vị và kiểm tra tính hợp lí của kết quả. Nhận định “Cường độ phóng xạ có thể tăng tùy ý bằng cách nung nóng mẫu.” sai vì trái với kiến thức trên.
}
\end{bt}

% ===== Câu 11 =====
% ID: L12C4B23-03-TLN
% Mức: TH
\begin{ex}
% ID: L12C4B23-03-TLN
% Mức: TH
Trong bốn phát biểu sau về nhận biết hiện tượng phóng xạ và các tia phóng xạ, có bao nhiêu phát biểu đúng? Chỉ ghi một số nguyên. (1) Cần đổi các đại lượng về hệ đơn vị thống nhất. (2) Phóng xạ là quá trình tự phát của hạt nhân không bền; tia $\alpha,\beta,\gamma$ có bản chất và khả năng đâm xuyên khác nhau. (3) Cường độ phóng xạ có thể tăng tùy ý bằng cách nung nóng mẫu. (4) Không cần kiểm tra dữ kiện.
\shortans{2}
\loigiai{
Đối chiếu từng phát biểu với kiến thức cốt lõi: Phóng xạ là quá trình tự phát của hạt nhân không bền; tia $\alpha,\beta,\gamma$ có bản chất và khả năng đâm xuyên khác nhau. Có 2 phát biểu đúng.
}
\end{ex}

% ===== Câu 12 =====
% ID: L12C4B23-03-TN
% Mức: NB
\begin{ex}
% ID: L12C4B23-03-TN
% Mức: NB
Chọn kết luận chính xác về nhận biết hiện tượng phóng xạ và các tia phóng xạ?
\choice
{Cường độ phóng xạ có thể tăng tùy ý bằng cách nung nóng mẫu.}
{Có thể bỏ qua điều kiện áp dụng và đơn vị trong mọi trường hợp.}
{\True Phóng xạ là quá trình tự phát của hạt nhân không bền; tia $\alpha,\beta,\gamma$ có bản chất và khả năng đâm xuyên khác nhau.}
{Kết quả của bài toán không phụ thuộc dữ kiện đã cho.}
\loigiai{
Phóng xạ là quá trình tự phát của hạt nhân không bền; tia $\alpha,\beta,\gamma$ có bản chất và khả năng đâm xuyên khác nhau. Vì vậy phương án $C$ đúng; các phương án còn lại trái với kiến thức hoặc điều kiện áp dụng.
}
\end{ex}

% ===== Câu 13 =====
% ID: L12C4B23-04-DS
% Mức: VD
\begin{ex}
% ID: L12C4B23-04-DS
% Mức: VD
Xét nhận biết hiện tượng phóng xạ và các tia phóng xạ (tình huống 4). Hãy xác định tính đúng, sai của bốn phát biểu sau.
\choiceTF
{Cường độ phóng xạ có thể tăng tùy ý bằng cách nung nóng mẫu.}
{\True Phóng xạ là quá trình tự phát của hạt nhân không bền; tia $\alpha,\beta,\gamma$ có bản chất và khả năng đâm xuyên khác nhau.}
{Có thể bỏ qua dữ kiện, đơn vị và ý nghĩa vật lí của kết quả.}
{\True Khi giải cần đọc đúng dữ kiện, dùng hệ đơn vị thống nhất và kiểm tra điều kiện áp dụng.}
\loigiai{
\begin{itemchoice}
\item (Sai) Cường độ phóng xạ có thể tăng tùy ý bằng cách nung nóng mẫu. (Vì): Cường độ phóng xạ có thể tăng tùy ý bằng cách nung nóng mẫu. b) Đúng: Phóng xạ là quá trình tự phát của hạt nhân không bền; tia $\alpha,\beta,\gamma$ có bản chất và khả năng đâm xuyên khác nhau. c) Sai: Có thể bỏ qua dữ kiện, đơn vị và ý nghĩa vật lí của kết quả. d) Đúng: Khi giải cần đọc đúng dữ kiện, dùng hệ đơn vị thống nhất và kiểm tra điều kiện áp dụng. Kiến thức cốt lõi: Phóng xạ là quá trình tự phát của hạt nhân không bền; tia $\alpha,\beta,\gamma$ có bản chất và khả năng đâm xuyên khác nhau
\item (Đúng) Phóng xạ là quá trình tự phát của hạt nhân không bền; tia $\alpha,\beta,\gamma$ có bản chất và khả năng đâm xuyên khác nhau.
\item (Sai) Có thể bỏ qua dữ kiện, đơn vị và ý nghĩa vật lí của kết quả.
\item (Đúng) Khi giải cần đọc đúng dữ kiện, dùng hệ đơn vị thống nhất và kiểm tra điều kiện áp dụng.
\end{itemchoice}
}
\end{ex}

% ===== Câu 14 =====
% ID: L12C4B23-04-TL
% Mức: VD
\begin{bt}
% ID: L12C4B23-04-TL
% Mức: VD
Đề xuất các bước giải một bài toán thuộc dạng nhận biết hiện tượng phóng xạ và các tia phóng xạ.
\loigiai{
Cần nêu được: Phóng xạ là quá trình tự phát của hạt nhân không bền; tia $\alpha,\beta,\gamma$ có bản chất và khả năng đâm xuyên khác nhau. Khi vận dụng phải xác định đúng dữ kiện, điều kiện áp dụng, hệ đơn vị và kiểm tra tính hợp lí của kết quả. Nhận định “Cường độ phóng xạ có thể tăng tùy ý bằng cách nung nóng mẫu.” sai vì trái với kiến thức trên.
}
\end{bt}

% ===== Câu 15 =====
% ID: L12C4B23-04-TLN
% Mức: VD
\begin{ex}
% ID: L12C4B23-04-TLN
% Mức: VD
Trong bốn phát biểu sau về nhận biết hiện tượng phóng xạ và các tia phóng xạ, có bao nhiêu phát biểu đúng? Chỉ ghi một số nguyên. (1) Phóng xạ là quá trình tự phát của hạt nhân không bền; tia $\alpha,\beta,\gamma$ có bản chất và khả năng đâm xuyên khác nhau. (2) Cần đối chiếu kết quả với điều kiện bài toán. (3) Cường độ phóng xạ có thể tăng tùy ý bằng cách nung nóng mẫu. (4) Mọi trường hợp đều cho cùng một kết quả.
\shortans{2}
\loigiai{
Đối chiếu từng phát biểu với kiến thức cốt lõi: Phóng xạ là quá trình tự phát của hạt nhân không bền; tia $\alpha,\beta,\gamma$ có bản chất và khả năng đâm xuyên khác nhau. Có 2 phát biểu đúng.
}
\end{ex}

% ===== Câu 16 =====
% ID: L12C4B23-04-TN
% Mức: TH
\begin{ex}
% ID: L12C4B23-04-TN
% Mức: TH
Nội dung nào mô tả đúng nhất về nhận biết hiện tượng phóng xạ và các tia phóng xạ?
\choice
{Cường độ phóng xạ có thể tăng tùy ý bằng cách nung nóng mẫu.}
{Có thể bỏ qua điều kiện áp dụng và đơn vị trong mọi trường hợp.}
{Kết quả của bài toán không phụ thuộc dữ kiện đã cho.}
{\True Phóng xạ là quá trình tự phát của hạt nhân không bền; tia $\alpha,\beta,\gamma$ có bản chất và khả năng đâm xuyên khác nhau.}
\loigiai{
Phóng xạ là quá trình tự phát của hạt nhân không bền; tia $\alpha,\beta,\gamma$ có bản chất và khả năng đâm xuyên khác nhau. Vì vậy phương án $D$ đúng; các phương án còn lại trái với kiến thức hoặc điều kiện áp dụng.
}
\end{ex}

% ===== Câu 17 =====
% ID: L12C4B23-05-DS
% Mức: VD
\begin{ex}
% ID: L12C4B23-05-DS
% Mức: VD
Xét nhận biết hiện tượng phóng xạ và các tia phóng xạ (tình huống 5). Hãy xác định tính đúng, sai của bốn phát biểu sau.
\choiceTF
{\True Phóng xạ là quá trình tự phát của hạt nhân không bền; tia $\alpha,\beta,\gamma$ có bản chất và khả năng đâm xuyên khác nhau.}
{Có thể bỏ qua dữ kiện, đơn vị và ý nghĩa vật lí của kết quả.}
{Cường độ phóng xạ có thể tăng tùy ý bằng cách nung nóng mẫu.}
{Có thể bỏ qua dữ kiện, đơn vị và ý nghĩa vật lí của kết quả.}
\loigiai{
\begin{itemchoice}
\item (Đúng) Phóng xạ là quá trình tự phát của hạt nhân không bền; tia $\alpha,\beta,\gamma$ có bản chất và khả năng đâm xuyên khác nhau. (Vì): Phóng xạ là quá trình tự phát của hạt nhân không bền; tia $\alpha,\beta,\gamma$ có bản chất và khả năng đâm xuyên khác nhau. b) Sai: Có thể bỏ qua dữ kiện, đơn vị và ý nghĩa vật lí của kết quả. c) Sai: Cường độ phóng xạ có thể tăng tùy ý bằng cách nung nóng mẫu. d) Sai: Có thể bỏ qua dữ kiện, đơn vị và ý nghĩa vật lí của kết quả. Kiến thức cốt lõi: Phóng xạ là quá trình tự phát của hạt nhân không bền; tia $\alpha,\beta,\gamma$ có bản chất và khả năng đâm xuyên khác nhau
\item (Sai) Có thể bỏ qua dữ kiện, đơn vị và ý nghĩa vật lí của kết quả.
\item (Sai) Cường độ phóng xạ có thể tăng tùy ý bằng cách nung nóng mẫu.
\item (Sai) Có thể bỏ qua dữ kiện, đơn vị và ý nghĩa vật lí của kết quả.
\end{itemchoice}
}
\end{ex}

% ===== Câu 18 =====
% ID: L12C4B23-05-TL
% Mức: VD
\begin{bt}
% ID: L12C4B23-05-TL
% Mức: VD
Nêu một tình huống thực tiễn liên quan đến nhận biết hiện tượng phóng xạ và các tia phóng xạ và giải thích bằng kiến thức Vật lí.
\loigiai{
Cần nêu được: Phóng xạ là quá trình tự phát của hạt nhân không bền; tia $\alpha,\beta,\gamma$ có bản chất và khả năng đâm xuyên khác nhau. Khi vận dụng phải xác định đúng dữ kiện, điều kiện áp dụng, hệ đơn vị và kiểm tra tính hợp lí của kết quả. Nhận định “Cường độ phóng xạ có thể tăng tùy ý bằng cách nung nóng mẫu.” sai vì trái với kiến thức trên.
}
\end{bt}

% ===== Câu 19 =====
% ID: L12C4B23-05-TLN
% Mức: VD
\begin{ex}
% ID: L12C4B23-05-TLN
% Mức: VD
Trong bốn phát biểu sau về nhận biết hiện tượng phóng xạ và các tia phóng xạ, có bao nhiêu phát biểu đúng? Chỉ ghi một số nguyên. (1) Cường độ phóng xạ có thể tăng tùy ý bằng cách nung nóng mẫu. (2) Cần đọc đúng dữ kiện và giả thiết. (3) Phóng xạ là quá trình tự phát của hạt nhân không bền; tia $\alpha,\beta,\gamma$ có bản chất và khả năng đâm xuyên khác nhau. (4) Có thể thay số trước khi chọn công thức.
\shortans{2}
\loigiai{
Đối chiếu từng phát biểu với kiến thức cốt lõi: Phóng xạ là quá trình tự phát của hạt nhân không bền; tia $\alpha,\beta,\gamma$ có bản chất và khả năng đâm xuyên khác nhau. Có 2 phát biểu đúng.
}
\end{ex}

% ===== Câu 20 =====
% ID: L12C4B23-05-TN
% Mức: TH
\begin{ex}
% ID: L12C4B23-05-TN
% Mức: TH
Khi phân tích, phát biểu nào đúng về nhận biết hiện tượng phóng xạ và các tia phóng xạ?
\choice
{\True Phóng xạ là quá trình tự phát của hạt nhân không bền; tia $\alpha,\beta,\gamma$ có bản chất và khả năng đâm xuyên khác nhau.}
{Cường độ phóng xạ có thể tăng tùy ý bằng cách nung nóng mẫu.}
{Có thể bỏ qua điều kiện áp dụng và đơn vị trong mọi trường hợp.}
{Kết quả của bài toán không phụ thuộc dữ kiện đã cho.}
\loigiai{
Phóng xạ là quá trình tự phát của hạt nhân không bền; tia $\alpha,\beta,\gamma$ có bản chất và khả năng đâm xuyên khác nhau. Vì vậy phương án $A$ đúng; các phương án còn lại trái với kiến thức hoặc điều kiện áp dụng.
}
\end{ex}

% ===== Câu 21 =====
% ID: L12C4B23-06-DS
% Mức: TH
\dangbt{Viết phương trình phóng xạ và xác định hạt nhân con}
\begin{ex}
% ID: L12C4B23-06-DS
% Mức: TH
Xét viết phương trình phóng xạ và xác định hạt nhân con (tình huống 1). Hãy xác định tính đúng, sai của bốn phát biểu sau.
\choiceTF
{\True Phóng xạ $\alpha$ làm $A$ giảm 4, $Z$ giảm 2; $\beta^-$ giữ nguyên $A$ và làm $Z$ tăng 1; $\gamma$ không đổi $A,Z$.}
{Phóng xạ $\gamma$ làm số khối giảm 1.}
{\True Khi giải cần đọc đúng dữ kiện, dùng hệ đơn vị thống nhất và kiểm tra điều kiện áp dụng.}
{Có thể bỏ qua dữ kiện, đơn vị và ý nghĩa vật lí của kết quả.}
\loigiai{
\begin{itemchoice}
\item (Đúng) Phóng xạ $\alpha$ làm $A$ giảm 4, $Z$ giảm 2; $\beta^-$ giữ nguyên $A$ và làm $Z$ tăng 1; $\gamma$ không đổi $A,Z$. (Vì): Phóng xạ $\alpha$ làm $A$ giảm 4, $Z$ giảm 2; $\beta^-$ giữ nguyên $A$ và làm $Z$ tăng 1; $\gamma$ không đổi $A,Z$. b) Sai: Phóng xạ $\gamma$ làm số khối giảm 1. c) Đúng: Khi giải cần đọc đúng dữ kiện, dùng hệ đơn vị thống nhất và kiểm tra điều kiện áp dụng. d) Sai: Có thể bỏ qua dữ kiện, đơn vị và ý nghĩa vật lí của kết quả. Kiến thức cốt lõi: Phóng xạ $\alpha$ làm $A$ giảm 4, $Z$ giảm 2; $\beta^-$ giữ nguyên $A$ và làm $Z$ tăng 1; $\gamma$ không đổi $A,Z$
\item (Sai) Phóng xạ $\gamma$ làm số khối giảm 1.
\item (Đúng) Khi giải cần đọc đúng dữ kiện, dùng hệ đơn vị thống nhất và kiểm tra điều kiện áp dụng.
\item (Sai) Có thể bỏ qua dữ kiện, đơn vị và ý nghĩa vật lí của kết quả.
\end{itemchoice}
}
\end{ex}

% ===== Câu 22 =====
% ID: L12C4B23-06-TL
% Mức: VDC
\dangbt{Nhận biết hiện tượng phóng xạ và các tia phóng xạ}
\begin{bt}
% ID: L12C4B23-06-TL
% Mức: VDC
So sánh nhận định đúng “Phóng xạ là quá trình tự phát của hạt nhân không bền; tia $\alpha,\beta,\gamma$ có bản chất và khả năng đâm xuyên khác nhau.” với nhận định sai “Cường độ phóng xạ có thể tăng tùy ý bằng cách nung nóng mẫu.”; chỉ rõ căn cứ Vật lí.
\loigiai{
Cần nêu được: Phóng xạ là quá trình tự phát của hạt nhân không bền; tia $\alpha,\beta,\gamma$ có bản chất và khả năng đâm xuyên khác nhau. Khi vận dụng phải xác định đúng dữ kiện, điều kiện áp dụng, hệ đơn vị và kiểm tra tính hợp lí của kết quả. Nhận định “Cường độ phóng xạ có thể tăng tùy ý bằng cách nung nóng mẫu.” sai vì trái với kiến thức trên.
}
\end{bt}

% ===== Câu 23 =====
% ID: L12C4B23-06-TLN
% Mức: VDC
\begin{ex}
% ID: L12C4B23-06-TLN
% Mức: VDC
Trong bốn phát biểu sau về nhận biết hiện tượng phóng xạ và các tia phóng xạ, có bao nhiêu phát biểu đúng? Chỉ ghi một số nguyên. (1) Kết quả cần có ý nghĩa vật lí. (2) Cường độ phóng xạ có thể tăng tùy ý bằng cách nung nóng mẫu. (3) Phải dùng đúng định luật cho tình huống. (4) Phóng xạ là quá trình tự phát của hạt nhân không bền; tia $\alpha,\beta,\gamma$ có bản chất và khả năng đâm xuyên khác nhau.
\shortans{3}
\loigiai{
Đối chiếu từng phát biểu với kiến thức cốt lõi: Phóng xạ là quá trình tự phát của hạt nhân không bền; tia $\alpha,\beta,\gamma$ có bản chất và khả năng đâm xuyên khác nhau. Có 3 phát biểu đúng.
}
\end{ex}

% ===== Câu 24 =====
% ID: L12C4B23-06-TN
% Mức: TH
\begin{ex}
% ID: L12C4B23-06-TN
% Mức: TH
Kết luận nào của học sinh là đúng về nhận biết hiện tượng phóng xạ và các tia phóng xạ?
\choice
{Cường độ phóng xạ có thể tăng tùy ý bằng cách nung nóng mẫu.}
{\True Phóng xạ là quá trình tự phát của hạt nhân không bền; tia $\alpha,\beta,\gamma$ có bản chất và khả năng đâm xuyên khác nhau.}
{Có thể bỏ qua điều kiện áp dụng và đơn vị trong mọi trường hợp.}
{Kết quả của bài toán không phụ thuộc dữ kiện đã cho.}
\loigiai{
Phóng xạ là quá trình tự phát của hạt nhân không bền; tia $\alpha,\beta,\gamma$ có bản chất và khả năng đâm xuyên khác nhau. Vì vậy phương án $B$ đúng; các phương án còn lại trái với kiến thức hoặc điều kiện áp dụng.
}
\end{ex}

% ===== Câu 25 =====
% ID: L12C4B23-07-DS
% Mức: TH
\dangbt{Viết phương trình phóng xạ và xác định hạt nhân con}
\begin{ex}
% ID: L12C4B23-07-DS
% Mức: TH
Xét viết phương trình phóng xạ và xác định hạt nhân con (tình huống 2). Hãy xác định tính đúng, sai của bốn phát biểu sau.
\choiceTF
{Phóng xạ $\gamma$ làm số khối giảm 1.}
{\True Phóng xạ $\alpha$ làm $A$ giảm 4, $Z$ giảm 2; $\beta^-$ giữ nguyên $A$ và làm $Z$ tăng 1; $\gamma$ không đổi $A,Z$.}
{Có thể bỏ qua dữ kiện, đơn vị và ý nghĩa vật lí của kết quả.}
{\True Khi giải cần đọc đúng dữ kiện, dùng hệ đơn vị thống nhất và kiểm tra điều kiện áp dụng.}
\loigiai{
\begin{itemchoice}
\item (Sai) Phóng xạ $\gamma$ làm số khối giảm 1. (Vì): Phóng xạ $\gamma$ làm số khối giảm 1. b) Đúng: Phóng xạ $\alpha$ làm $A$ giảm 4, $Z$ giảm 2; $\beta^-$ giữ nguyên $A$ và làm $Z$ tăng 1; $\gamma$ không đổi $A,Z$. c) Sai: Có thể bỏ qua dữ kiện, đơn vị và ý nghĩa vật lí của kết quả. d) Đúng: Khi giải cần đọc đúng dữ kiện, dùng hệ đơn vị thống nhất và kiểm tra điều kiện áp dụng. Kiến thức cốt lõi: Phóng xạ $\alpha$ làm $A$ giảm 4, $Z$ giảm 2; $\beta^-$ giữ nguyên $A$ và làm $Z$ tăng 1; $\gamma$ không đổi $A,Z$
\item (Đúng) Phóng xạ $\alpha$ làm $A$ giảm 4, $Z$ giảm 2; $\beta^-$ giữ nguyên $A$ và làm $Z$ tăng 1; $\gamma$ không đổi $A,Z$.
\item (Sai) Có thể bỏ qua dữ kiện, đơn vị và ý nghĩa vật lí của kết quả.
\item (Đúng) Khi giải cần đọc đúng dữ kiện, dùng hệ đơn vị thống nhất và kiểm tra điều kiện áp dụng.
\end{itemchoice}
}
\end{ex}

% ===== Câu 26 =====
% ID: L12C4B23-07-TL
% Mức: TH
\begin{bt}
% ID: L12C4B23-07-TL
% Mức: TH
Trình bày kiến thức cốt lõi của viết phương trình phóng xạ và xác định hạt nhân con và nêu điều kiện áp dụng.
\loigiai{
Cần nêu được: Phóng xạ $\alpha$ làm $A$ giảm 4, $Z$ giảm 2; $\beta^-$ giữ nguyên $A$ và làm $Z$ tăng 1; $\gamma$ không đổi $A,Z$. Khi vận dụng phải xác định đúng dữ kiện, điều kiện áp dụng, hệ đơn vị và kiểm tra tính hợp lí của kết quả. Nhận định “Phóng xạ $\gamma$ làm số khối giảm 1.” sai vì trái với kiến thức trên.
}
\end{bt}

% ===== Câu 27 =====
% ID: L12C4B23-07-TLN
% Mức: TH
\begin{ex}
% ID: L12C4B23-07-TLN
% Mức: TH
Trong bốn phát biểu sau về viết phương trình phóng xạ và xác định hạt nhân con, có bao nhiêu phát biểu đúng? Chỉ ghi một số nguyên. (1) Phóng xạ $\alpha$ làm $A$ giảm 4, $Z$ giảm 2; $\beta^-$ giữ nguyên $A$ và làm $Z$ tăng 1; $\gamma$ không đổi $A,Z$. (2) Phóng xạ $\gamma$ làm số khối giảm 1. (3) Cần kiểm tra điều kiện áp dụng trước khi tính toán. (4) Có thể bỏ qua đơn vị của kết quả.
\shortans{2}
\loigiai{
Đối chiếu từng phát biểu với kiến thức cốt lõi: Phóng xạ $\alpha$ làm $A$ giảm 4, $Z$ giảm 2; $\beta^-$ giữ nguyên $A$ và làm $Z$ tăng 1; $\gamma$ không đổi $A,Z$. Có 2 phát biểu đúng.
}
\end{ex}

% ===== Câu 28 =====
% ID: L12C4B23-07-TN
% Mức: TH
\dangbt{Nhận biết hiện tượng phóng xạ và các tia phóng xạ}
\begin{ex}
% ID: L12C4B23-07-TN
% Mức: TH
Chọn phát biểu không trái với kiến thức về nhận biết hiện tượng phóng xạ và các tia phóng xạ?
\choice
{Cường độ phóng xạ có thể tăng tùy ý bằng cách nung nóng mẫu.}
{Có thể bỏ qua điều kiện áp dụng và đơn vị trong mọi trường hợp.}
{\True Phóng xạ là quá trình tự phát của hạt nhân không bền; tia $\alpha,\beta,\gamma$ có bản chất và khả năng đâm xuyên khác nhau.}
{Kết quả của bài toán không phụ thuộc dữ kiện đã cho.}
\loigiai{
Phóng xạ là quá trình tự phát của hạt nhân không bền; tia $\alpha,\beta,\gamma$ có bản chất và khả năng đâm xuyên khác nhau. Vì vậy phương án $C$ đúng; các phương án còn lại trái với kiến thức hoặc điều kiện áp dụng.
}
\end{ex}

% ===== Câu 29 =====
% ID: L12C4B23-08-DS
% Mức: VD
\dangbt{Viết phương trình phóng xạ và xác định hạt nhân con}
\begin{ex}
% ID: L12C4B23-08-DS
% Mức: VD
Xét viết phương trình phóng xạ và xác định hạt nhân con (tình huống 3). Hãy xác định tính đúng, sai của bốn phát biểu sau.
\choiceTF
{\True Phóng xạ $\alpha$ làm $A$ giảm 4, $Z$ giảm 2; $\beta^-$ giữ nguyên $A$ và làm $Z$ tăng 1; $\gamma$ không đổi $A,Z$.}
{Phóng xạ $\gamma$ làm số khối giảm 1.}
{\True Khi giải cần đọc đúng dữ kiện, dùng hệ đơn vị thống nhất và kiểm tra điều kiện áp dụng.}
{Có thể bỏ qua dữ kiện, đơn vị và ý nghĩa vật lí của kết quả.}
\loigiai{
\begin{itemchoice}
\item (Đúng) Phóng xạ $\alpha$ làm $A$ giảm 4, $Z$ giảm 2; $\beta^-$ giữ nguyên $A$ và làm $Z$ tăng 1; $\gamma$ không đổi $A,Z$. (Vì): Phóng xạ $\alpha$ làm $A$ giảm 4, $Z$ giảm 2; $\beta^-$ giữ nguyên $A$ và làm $Z$ tăng 1; $\gamma$ không đổi $A,Z$. b) Sai: Phóng xạ $\gamma$ làm số khối giảm 1. c) Đúng: Khi giải cần đọc đúng dữ kiện, dùng hệ đơn vị thống nhất và kiểm tra điều kiện áp dụng. d) Sai: Có thể bỏ qua dữ kiện, đơn vị và ý nghĩa vật lí của kết quả. Kiến thức cốt lõi: Phóng xạ $\alpha$ làm $A$ giảm 4, $Z$ giảm 2; $\beta^-$ giữ nguyên $A$ và làm $Z$ tăng 1; $\gamma$ không đổi $A,Z$
\item (Sai) Phóng xạ $\gamma$ làm số khối giảm 1.
\item (Đúng) Khi giải cần đọc đúng dữ kiện, dùng hệ đơn vị thống nhất và kiểm tra điều kiện áp dụng.
\item (Sai) Có thể bỏ qua dữ kiện, đơn vị và ý nghĩa vật lí của kết quả.
\end{itemchoice}
}
\end{ex}

% ===== Câu 30 =====
% ID: L12C4B23-08-TL
% Mức: TH
\begin{bt}
% ID: L12C4B23-08-TL
% Mức: TH
Giải thích vì sao nhận định sau đúng: “Phóng xạ $\alpha$ làm $A$ giảm 4, $Z$ giảm 2; $\beta^-$ giữ nguyên $A$ và làm $Z$ tăng 1; $\gamma$ không đổi $A,Z$.”
\loigiai{
Cần nêu được: Phóng xạ $\alpha$ làm $A$ giảm 4, $Z$ giảm 2; $\beta^-$ giữ nguyên $A$ và làm $Z$ tăng 1; $\gamma$ không đổi $A,Z$. Khi vận dụng phải xác định đúng dữ kiện, điều kiện áp dụng, hệ đơn vị và kiểm tra tính hợp lí của kết quả. Nhận định “Phóng xạ $\gamma$ làm số khối giảm 1.” sai vì trái với kiến thức trên.
}
\end{bt}

% ===== Câu 31 =====
% ID: L12C4B23-08-TLN
% Mức: TH
\begin{ex}
% ID: L12C4B23-08-TLN
% Mức: TH
Trong bốn phát biểu sau về viết phương trình phóng xạ và xác định hạt nhân con, có bao nhiêu phát biểu đúng? Chỉ ghi một số nguyên. (1) Phóng xạ $\gamma$ làm số khối giảm 1. (2) Phóng xạ $\alpha$ làm $A$ giảm 4, $Z$ giảm 2; $\beta^-$ giữ nguyên $A$ và làm $Z$ tăng 1; $\gamma$ không đổi $A,Z$. (3) Kết quả phải phù hợp ý nghĩa vật lí. (4) Mọi đại lượng đều có cùng bản chất.
\shortans{2}
\loigiai{
Đối chiếu từng phát biểu với kiến thức cốt lõi: Phóng xạ $\alpha$ làm $A$ giảm 4, $Z$ giảm 2; $\beta^-$ giữ nguyên $A$ và làm $Z$ tăng 1; $\gamma$ không đổi $A,Z$. Có 2 phát biểu đúng.
}
\end{ex}

% ===== Câu 32 =====
% ID: L12C4B23-08-TN
% Mức: VD
\dangbt{Nhận biết hiện tượng phóng xạ và các tia phóng xạ}
\begin{ex}
% ID: L12C4B23-08-TN
% Mức: VD
Cơ sở Vật lí đúng để giải bài là gì về nhận biết hiện tượng phóng xạ và các tia phóng xạ?
\choice
{Cường độ phóng xạ có thể tăng tùy ý bằng cách nung nóng mẫu.}
{Có thể bỏ qua điều kiện áp dụng và đơn vị trong mọi trường hợp.}
{Kết quả của bài toán không phụ thuộc dữ kiện đã cho.}
{\True Phóng xạ là quá trình tự phát của hạt nhân không bền; tia $\alpha,\beta,\gamma$ có bản chất và khả năng đâm xuyên khác nhau.}
\loigiai{
Phóng xạ là quá trình tự phát của hạt nhân không bền; tia $\alpha,\beta,\gamma$ có bản chất và khả năng đâm xuyên khác nhau. Vì vậy phương án $D$ đúng; các phương án còn lại trái với kiến thức hoặc điều kiện áp dụng.
}
\end{ex}

% ===== Câu 33 =====
% ID: L12C4B23-09-DS
% Mức: VD
\dangbt{Viết phương trình phóng xạ và xác định hạt nhân con}
\begin{ex}
% ID: L12C4B23-09-DS
% Mức: VD
Xét viết phương trình phóng xạ và xác định hạt nhân con (tình huống 4). Hãy xác định tính đúng, sai của bốn phát biểu sau.
\choiceTF
{Phóng xạ $\gamma$ làm số khối giảm 1.}
{\True Phóng xạ $\alpha$ làm $A$ giảm 4, $Z$ giảm 2; $\beta^-$ giữ nguyên $A$ và làm $Z$ tăng 1; $\gamma$ không đổi $A,Z$.}
{Có thể bỏ qua dữ kiện, đơn vị và ý nghĩa vật lí của kết quả.}
{\True Khi giải cần đọc đúng dữ kiện, dùng hệ đơn vị thống nhất và kiểm tra điều kiện áp dụng.}
\loigiai{
\begin{itemchoice}
\item (Sai) Phóng xạ $\gamma$ làm số khối giảm 1. (Vì): Phóng xạ $\gamma$ làm số khối giảm 1. b) Đúng: Phóng xạ $\alpha$ làm $A$ giảm 4, $Z$ giảm 2; $\beta^-$ giữ nguyên $A$ và làm $Z$ tăng 1; $\gamma$ không đổi $A,Z$. c) Sai: Có thể bỏ qua dữ kiện, đơn vị và ý nghĩa vật lí của kết quả. d) Đúng: Khi giải cần đọc đúng dữ kiện, dùng hệ đơn vị thống nhất và kiểm tra điều kiện áp dụng. Kiến thức cốt lõi: Phóng xạ $\alpha$ làm $A$ giảm 4, $Z$ giảm 2; $\beta^-$ giữ nguyên $A$ và làm $Z$ tăng 1; $\gamma$ không đổi $A,Z$
\item (Đúng) Phóng xạ $\alpha$ làm $A$ giảm 4, $Z$ giảm 2; $\beta^-$ giữ nguyên $A$ và làm $Z$ tăng 1; $\gamma$ không đổi $A,Z$.
\item (Sai) Có thể bỏ qua dữ kiện, đơn vị và ý nghĩa vật lí của kết quả.
\item (Đúng) Khi giải cần đọc đúng dữ kiện, dùng hệ đơn vị thống nhất và kiểm tra điều kiện áp dụng.
\end{itemchoice}
}
\end{ex}

% ===== Câu 34 =====
% ID: L12C4B23-09-TL
% Mức: VD
\begin{bt}
% ID: L12C4B23-09-TL
% Mức: VD
Phân tích sai lầm trong nhận định: “Phóng xạ $\gamma$ làm số khối giảm 1.”
\loigiai{
Cần nêu được: Phóng xạ $\alpha$ làm $A$ giảm 4, $Z$ giảm 2; $\beta^-$ giữ nguyên $A$ và làm $Z$ tăng 1; $\gamma$ không đổi $A,Z$. Khi vận dụng phải xác định đúng dữ kiện, điều kiện áp dụng, hệ đơn vị và kiểm tra tính hợp lí của kết quả. Nhận định “Phóng xạ $\gamma$ làm số khối giảm 1.” sai vì trái với kiến thức trên.
}
\end{bt}

% ===== Câu 35 =====
% ID: L12C4B23-09-TLN
% Mức: TH
\begin{ex}
% ID: L12C4B23-09-TLN
% Mức: TH
Trong bốn phát biểu sau về viết phương trình phóng xạ và xác định hạt nhân con, có bao nhiêu phát biểu đúng? Chỉ ghi một số nguyên. (1) Cần đổi các đại lượng về hệ đơn vị thống nhất. (2) Phóng xạ $\alpha$ làm $A$ giảm 4, $Z$ giảm 2; $\beta^-$ giữ nguyên $A$ và làm $Z$ tăng 1; $\gamma$ không đổi $A,Z$. (3) Phóng xạ $\gamma$ làm số khối giảm 1. (4) Không cần kiểm tra dữ kiện.
\shortans{2}
\loigiai{
Đối chiếu từng phát biểu với kiến thức cốt lõi: Phóng xạ $\alpha$ làm $A$ giảm 4, $Z$ giảm 2; $\beta^-$ giữ nguyên $A$ và làm $Z$ tăng 1; $\gamma$ không đổi $A,Z$. Có 2 phát biểu đúng.
}
\end{ex}

% ===== Câu 36 =====
% ID: L12C4B23-09-TN
% Mức: VD
\dangbt{Nhận biết hiện tượng phóng xạ và các tia phóng xạ}
\begin{ex}
% ID: L12C4B23-09-TN
% Mức: VD
Trong một bài thực hành, nhận xét nào đúng về nhận biết hiện tượng phóng xạ và các tia phóng xạ?
\choice
{\True Phóng xạ là quá trình tự phát của hạt nhân không bền; tia $\alpha,\beta,\gamma$ có bản chất và khả năng đâm xuyên khác nhau.}
{Cường độ phóng xạ có thể tăng tùy ý bằng cách nung nóng mẫu.}
{Có thể bỏ qua điều kiện áp dụng và đơn vị trong mọi trường hợp.}
{Kết quả của bài toán không phụ thuộc dữ kiện đã cho.}
\loigiai{
Phóng xạ là quá trình tự phát của hạt nhân không bền; tia $\alpha,\beta,\gamma$ có bản chất và khả năng đâm xuyên khác nhau. Vì vậy phương án $A$ đúng; các phương án còn lại trái với kiến thức hoặc điều kiện áp dụng.
}
\end{ex}

% ===== Câu 37 =====
% ID: L12C4B23-10-DS
% Mức: VD
\dangbt{Viết phương trình phóng xạ và xác định hạt nhân con}
\begin{ex}
% ID: L12C4B23-10-DS
% Mức: VD
Xét viết phương trình phóng xạ và xác định hạt nhân con (tình huống 5). Hãy xác định tính đúng, sai của bốn phát biểu sau.
\choiceTF
{\True Phóng xạ $\alpha$ làm $A$ giảm 4, $Z$ giảm 2; $\beta^-$ giữ nguyên $A$ và làm $Z$ tăng 1; $\gamma$ không đổi $A,Z$.}
{Có thể bỏ qua dữ kiện, đơn vị và ý nghĩa vật lí của kết quả.}
{Phóng xạ $\gamma$ làm số khối giảm 1.}
{Có thể bỏ qua dữ kiện, đơn vị và ý nghĩa vật lí của kết quả.}
\loigiai{
\begin{itemchoice}
\item (Đúng) Phóng xạ $\alpha$ làm $A$ giảm 4, $Z$ giảm 2; $\beta^-$ giữ nguyên $A$ và làm $Z$ tăng 1; $\gamma$ không đổi $A,Z$. (Vì): Phóng xạ $\alpha$ làm $A$ giảm 4, $Z$ giảm 2; $\beta^-$ giữ nguyên $A$ và làm $Z$ tăng 1; $\gamma$ không đổi $A,Z$. b) Sai: Có thể bỏ qua dữ kiện, đơn vị và ý nghĩa vật lí của kết quả. c) Sai: Phóng xạ $\gamma$ làm số khối giảm 1. d) Sai: Có thể bỏ qua dữ kiện, đơn vị và ý nghĩa vật lí của kết quả. Kiến thức cốt lõi: Phóng xạ $\alpha$ làm $A$ giảm 4, $Z$ giảm 2; $\beta^-$ giữ nguyên $A$ và làm $Z$ tăng 1; $\gamma$ không đổi $A,Z$
\item (Sai) Có thể bỏ qua dữ kiện, đơn vị và ý nghĩa vật lí của kết quả.
\item (Sai) Phóng xạ $\gamma$ làm số khối giảm 1.
\item (Sai) Có thể bỏ qua dữ kiện, đơn vị và ý nghĩa vật lí của kết quả.
\end{itemchoice}
}
\end{ex}

% ===== Câu 38 =====
% ID: L12C4B23-10-TL
% Mức: VD
\begin{bt}
% ID: L12C4B23-10-TL
% Mức: VD
Đề xuất các bước giải một bài toán thuộc dạng viết phương trình phóng xạ và xác định hạt nhân con.
\loigiai{
Cần nêu được: Phóng xạ $\alpha$ làm $A$ giảm 4, $Z$ giảm 2; $\beta^-$ giữ nguyên $A$ và làm $Z$ tăng 1; $\gamma$ không đổi $A,Z$. Khi vận dụng phải xác định đúng dữ kiện, điều kiện áp dụng, hệ đơn vị và kiểm tra tính hợp lí của kết quả. Nhận định “Phóng xạ $\gamma$ làm số khối giảm 1.” sai vì trái với kiến thức trên.
}
\end{bt}

% ===== Câu 39 =====
% ID: L12C4B23-10-TLN
% Mức: VD
\begin{ex}
% ID: L12C4B23-10-TLN
% Mức: VD
Trong bốn phát biểu sau về viết phương trình phóng xạ và xác định hạt nhân con, có bao nhiêu phát biểu đúng? Chỉ ghi một số nguyên. (1) Phóng xạ $\alpha$ làm $A$ giảm 4, $Z$ giảm 2; $\beta^-$ giữ nguyên $A$ và làm $Z$ tăng 1; $\gamma$ không đổi $A,Z$. (2) Cần đối chiếu kết quả với điều kiện bài toán. (3) Phóng xạ $\gamma$ làm số khối giảm 1. (4) Mọi trường hợp đều cho cùng một kết quả.
\shortans{2}
\loigiai{
Đối chiếu từng phát biểu với kiến thức cốt lõi: Phóng xạ $\alpha$ làm $A$ giảm 4, $Z$ giảm 2; $\beta^-$ giữ nguyên $A$ và làm $Z$ tăng 1; $\gamma$ không đổi $A,Z$. Có 2 phát biểu đúng.
}
\end{ex}

% ===== Câu 40 =====
% ID: L12C4B23-10-TN
% Mức: VD
\dangbt{Nhận biết hiện tượng phóng xạ và các tia phóng xạ}
\begin{ex}
% ID: L12C4B23-10-TN
% Mức: VD
Kết luận đúng khi vận dụng là về nhận biết hiện tượng phóng xạ và các tia phóng xạ?
\choice
{Cường độ phóng xạ có thể tăng tùy ý bằng cách nung nóng mẫu.}
{\True Phóng xạ là quá trình tự phát của hạt nhân không bền; tia $\alpha,\beta,\gamma$ có bản chất và khả năng đâm xuyên khác nhau.}
{Có thể bỏ qua điều kiện áp dụng và đơn vị trong mọi trường hợp.}
{Kết quả của bài toán không phụ thuộc dữ kiện đã cho.}
\loigiai{
Phóng xạ là quá trình tự phát của hạt nhân không bền; tia $\alpha,\beta,\gamma$ có bản chất và khả năng đâm xuyên khác nhau. Vì vậy phương án $B$ đúng; các phương án còn lại trái với kiến thức hoặc điều kiện áp dụng.
}
\end{ex}

% ===== Câu 41 =====
% ID: L12C4B23-11-DS
% Mức: TH
\dangbt{Vận dụng định luật phóng xạ theo số hạt và khối lượng}
\begin{ex}
% ID: L12C4B23-11-DS
% Mức: TH
Xét vận dụng định luật phóng xạ theo số hạt và khối lượng (tình huống 1). Hãy xác định tính đúng, sai của bốn phát biểu sau.
\choiceTF
{\True Số hạt chưa phân rã và khối lượng chất phóng xạ giảm theo $N=N_0 2^{-t/T}$ và $m=m_0 2^{-t/T}$.}
{Sau mỗi chu kì bán rã, lượng chất chưa phân rã giảm một lượng cố định như nhau.}
{\True Khi giải cần đọc đúng dữ kiện, dùng hệ đơn vị thống nhất và kiểm tra điều kiện áp dụng.}
{Có thể bỏ qua dữ kiện, đơn vị và ý nghĩa vật lí của kết quả.}
\loigiai{
\begin{itemchoice}
\item (Đúng) Số hạt chưa phân rã và khối lượng chất phóng xạ giảm theo $N=N_0 2^{-t/T}$ và $m=m_0 2^{-t/T}$. (Vì): Số hạt chưa phân rã và khối lượng chất phóng xạ giảm theo $N=N_0 2^{-t/T}$ và $m=m_0 2^{-t/T}$. b) Sai: Sau mỗi chu kì bán rã, lượng chất chưa phân rã giảm một lượng cố định như nhau. c) Đúng: Khi giải cần đọc đúng dữ kiện, dùng hệ đơn vị thống nhất và kiểm tra điều kiện áp dụng. d) Sai: Có thể bỏ qua dữ kiện, đơn vị và ý nghĩa vật lí của kết quả. Kiến thức cốt lõi: Số hạt chưa phân rã và khối lượng chất phóng xạ giảm theo $N=N_0 2^{-t/T}$ và $m=m_0 2^{-t/T}$
\item (Sai) Sau mỗi chu kì bán rã, lượng chất chưa phân rã giảm một lượng cố định như nhau.
\item (Đúng) Khi giải cần đọc đúng dữ kiện, dùng hệ đơn vị thống nhất và kiểm tra điều kiện áp dụng.
\item (Sai) Có thể bỏ qua dữ kiện, đơn vị và ý nghĩa vật lí của kết quả.
\end{itemchoice}
}
\end{ex}

% ===== Câu 42 =====
% ID: L12C4B23-11-TL
% Mức: VD
\dangbt{Viết phương trình phóng xạ và xác định hạt nhân con}
\begin{bt}
% ID: L12C4B23-11-TL
% Mức: VD
Nêu một tình huống thực tiễn liên quan đến viết phương trình phóng xạ và xác định hạt nhân con và giải thích bằng kiến thức Vật lí.
\loigiai{
Cần nêu được: Phóng xạ $\alpha$ làm $A$ giảm 4, $Z$ giảm 2; $\beta^-$ giữ nguyên $A$ và làm $Z$ tăng 1; $\gamma$ không đổi $A,Z$. Khi vận dụng phải xác định đúng dữ kiện, điều kiện áp dụng, hệ đơn vị và kiểm tra tính hợp lí của kết quả. Nhận định “Phóng xạ $\gamma$ làm số khối giảm 1.” sai vì trái với kiến thức trên.
}
\end{bt}

% ===== Câu 43 =====
% ID: L12C4B23-11-TLN
% Mức: VD
\begin{ex}
% ID: L12C4B23-11-TLN
% Mức: VD
Trong bốn phát biểu sau về viết phương trình phóng xạ và xác định hạt nhân con, có bao nhiêu phát biểu đúng? Chỉ ghi một số nguyên. (1) Phóng xạ $\gamma$ làm số khối giảm 1. (2) Cần đọc đúng dữ kiện và giả thiết. (3) Phóng xạ $\alpha$ làm $A$ giảm 4, $Z$ giảm 2; $\beta^-$ giữ nguyên $A$ và làm $Z$ tăng 1; $\gamma$ không đổi $A,Z$. (4) Có thể thay số trước khi chọn công thức.
\shortans{2}
\loigiai{
Đối chiếu từng phát biểu với kiến thức cốt lõi: Phóng xạ $\alpha$ làm $A$ giảm 4, $Z$ giảm 2; $\beta^-$ giữ nguyên $A$ và làm $Z$ tăng 1; $\gamma$ không đổi $A,Z$. Có 2 phát biểu đúng.
}
\end{ex}

% ===== Câu 44 =====
% ID: L12C4B23-11-TN
% Mức: NB
\begin{ex}
% ID: L12C4B23-11-TN
% Mức: NB
Phát biểu nào sau đây đúng về viết phương trình phóng xạ và xác định hạt nhân con?
\choice
{\True Phóng xạ $\alpha$ làm $A$ giảm 4, $Z$ giảm 2; $\beta^-$ giữ nguyên $A$ và làm $Z$ tăng 1; $\gamma$ không đổi $A,Z$.}
{Phóng xạ $\gamma$ làm số khối giảm 1.}
{Có thể bỏ qua điều kiện áp dụng và đơn vị trong mọi trường hợp.}
{Kết quả của bài toán không phụ thuộc dữ kiện đã cho.}
\loigiai{
Phóng xạ $\alpha$ làm $A$ giảm 4, $Z$ giảm 2; $\beta^-$ giữ nguyên $A$ và làm $Z$ tăng 1; $\gamma$ không đổi $A,Z$. Vì vậy phương án $A$ đúng; các phương án còn lại trái với kiến thức hoặc điều kiện áp dụng.
}
\end{ex}

% ===== Câu 45 =====
% ID: L12C4B23-12-DS
% Mức: TH
\dangbt{Vận dụng định luật phóng xạ theo số hạt và khối lượng}
\begin{ex}
% ID: L12C4B23-12-DS
% Mức: TH
Xét vận dụng định luật phóng xạ theo số hạt và khối lượng (tình huống 2). Hãy xác định tính đúng, sai của bốn phát biểu sau.
\choiceTF
{Sau mỗi chu kì bán rã, lượng chất chưa phân rã giảm một lượng cố định như nhau.}
{\True Số hạt chưa phân rã và khối lượng chất phóng xạ giảm theo $N=N_0 2^{-t/T}$ và $m=m_0 2^{-t/T}$.}
{Có thể bỏ qua dữ kiện, đơn vị và ý nghĩa vật lí của kết quả.}
{\True Khi giải cần đọc đúng dữ kiện, dùng hệ đơn vị thống nhất và kiểm tra điều kiện áp dụng.}
\loigiai{
\begin{itemchoice}
\item (Sai) Sau mỗi chu kì bán rã, lượng chất chưa phân rã giảm một lượng cố định như nhau. (Vì): au mỗi chu kì bán rã, lượng chất chưa phân rã giảm một lượng cố định như nhau. b) Đúng: Số hạt chưa phân rã và khối lượng chất phóng xạ giảm theo $N=N_0 2^{-t/T}$ và $m=m_0 2^{-t/T}$. c) Sai: Có thể bỏ qua dữ kiện, đơn vị và ý nghĩa vật lí của kết quả. d) Đúng: Khi giải cần đọc đúng dữ kiện, dùng hệ đơn vị thống nhất và kiểm tra điều kiện áp dụng. Kiến thức cốt lõi: Số hạt chưa phân rã và khối lượng chất phóng xạ giảm theo $N=N_0 2^{-t/T}$ và $m=m_0 2^{-t/T}$
\item (Đúng) Số hạt chưa phân rã và khối lượng chất phóng xạ giảm theo $N=N_0 2^{-t/T}$ và $m=m_0 2^{-t/T}$.
\item (Sai) Có thể bỏ qua dữ kiện, đơn vị và ý nghĩa vật lí của kết quả.
\item (Đúng) Khi giải cần đọc đúng dữ kiện, dùng hệ đơn vị thống nhất và kiểm tra điều kiện áp dụng.
\end{itemchoice}
}
\end{ex}

% ===== Câu 46 =====
% ID: L12C4B23-12-TL
% Mức: VDC
\dangbt{Viết phương trình phóng xạ và xác định hạt nhân con}
\begin{bt}
% ID: L12C4B23-12-TL
% Mức: VDC
So sánh nhận định đúng “Phóng xạ $\alpha$ làm $A$ giảm 4, $Z$ giảm 2; $\beta^-$ giữ nguyên $A$ và làm $Z$ tăng 1; $\gamma$ không đổi $A,Z$.” với nhận định sai “Phóng xạ $\gamma$ làm số khối giảm 1.”; chỉ rõ căn cứ Vật lí.
\loigiai{
Cần nêu được: Phóng xạ $\alpha$ làm $A$ giảm 4, $Z$ giảm 2; $\beta^-$ giữ nguyên $A$ và làm $Z$ tăng 1; $\gamma$ không đổi $A,Z$. Khi vận dụng phải xác định đúng dữ kiện, điều kiện áp dụng, hệ đơn vị và kiểm tra tính hợp lí của kết quả. Nhận định “Phóng xạ $\gamma$ làm số khối giảm 1.” sai vì trái với kiến thức trên.
}
\end{bt}

% ===== Câu 47 =====
% ID: L12C4B23-12-TLN
% Mức: VDC
\begin{ex}
% ID: L12C4B23-12-TLN
% Mức: VDC
Trong bốn phát biểu sau về viết phương trình phóng xạ và xác định hạt nhân con, có bao nhiêu phát biểu đúng? Chỉ ghi một số nguyên. (1) Kết quả cần có ý nghĩa vật lí. (2) Phóng xạ $\gamma$ làm số khối giảm 1. (3) Phải dùng đúng định luật cho tình huống. (4) Phóng xạ $\alpha$ làm $A$ giảm 4, $Z$ giảm 2; $\beta^-$ giữ nguyên $A$ và làm $Z$ tăng 1; $\gamma$ không đổi $A,Z$.
\shortans{3}
\loigiai{
Đối chiếu từng phát biểu với kiến thức cốt lõi: Phóng xạ $\alpha$ làm $A$ giảm 4, $Z$ giảm 2; $\beta^-$ giữ nguyên $A$ và làm $Z$ tăng 1; $\gamma$ không đổi $A,Z$. Có 3 phát biểu đúng.
}
\end{ex}

% ===== Câu 48 =====
% ID: L12C4B23-12-TN
% Mức: NB
\begin{ex}
% ID: L12C4B23-12-TN
% Mức: NB
Nhận định nào phù hợp với kiến thức Vật lí về viết phương trình phóng xạ và xác định hạt nhân con?
\choice
{Phóng xạ $\gamma$ làm số khối giảm 1.}
{\True Phóng xạ $\alpha$ làm $A$ giảm 4, $Z$ giảm 2; $\beta^-$ giữ nguyên $A$ và làm $Z$ tăng 1; $\gamma$ không đổi $A,Z$.}
{Có thể bỏ qua điều kiện áp dụng và đơn vị trong mọi trường hợp.}
{Kết quả của bài toán không phụ thuộc dữ kiện đã cho.}
\loigiai{
Phóng xạ $\alpha$ làm $A$ giảm 4, $Z$ giảm 2; $\beta^-$ giữ nguyên $A$ và làm $Z$ tăng 1; $\gamma$ không đổi $A,Z$. Vì vậy phương án $B$ đúng; các phương án còn lại trái với kiến thức hoặc điều kiện áp dụng.
}
\end{ex}

% ===== Câu 49 =====
% ID: L12C4B23-13-DS
% Mức: VD
\dangbt{Vận dụng định luật phóng xạ theo số hạt và khối lượng}
\begin{ex}
% ID: L12C4B23-13-DS
% Mức: VD
Xét vận dụng định luật phóng xạ theo số hạt và khối lượng (tình huống 3). Hãy xác định tính đúng, sai của bốn phát biểu sau.
\choiceTF
{\True Số hạt chưa phân rã và khối lượng chất phóng xạ giảm theo $N=N_0 2^{-t/T}$ và $m=m_0 2^{-t/T}$.}
{Sau mỗi chu kì bán rã, lượng chất chưa phân rã giảm một lượng cố định như nhau.}
{\True Khi giải cần đọc đúng dữ kiện, dùng hệ đơn vị thống nhất và kiểm tra điều kiện áp dụng.}
{Có thể bỏ qua dữ kiện, đơn vị và ý nghĩa vật lí của kết quả.}
\loigiai{
\begin{itemchoice}
\item (Đúng) Số hạt chưa phân rã và khối lượng chất phóng xạ giảm theo $N=N_0 2^{-t/T}$ và $m=m_0 2^{-t/T}$. (Vì): Số hạt chưa phân rã và khối lượng chất phóng xạ giảm theo $N=N_0 2^{-t/T}$ và $m=m_0 2^{-t/T}$. b) Sai: Sau mỗi chu kì bán rã, lượng chất chưa phân rã giảm một lượng cố định như nhau. c) Đúng: Khi giải cần đọc đúng dữ kiện, dùng hệ đơn vị thống nhất và kiểm tra điều kiện áp dụng. d) Sai: Có thể bỏ qua dữ kiện, đơn vị và ý nghĩa vật lí của kết quả. Kiến thức cốt lõi: Số hạt chưa phân rã và khối lượng chất phóng xạ giảm theo $N=N_0 2^{-t/T}$ và $m=m_0 2^{-t/T}$
\item (Sai) Sau mỗi chu kì bán rã, lượng chất chưa phân rã giảm một lượng cố định như nhau.
\item (Đúng) Khi giải cần đọc đúng dữ kiện, dùng hệ đơn vị thống nhất và kiểm tra điều kiện áp dụng.
\item (Sai) Có thể bỏ qua dữ kiện, đơn vị và ý nghĩa vật lí của kết quả.
\end{itemchoice}
}
\end{ex}

% ===== Câu 50 =====
% ID: L12C4B23-13-TL
% Mức: TH
\begin{bt}
% ID: L12C4B23-13-TL
% Mức: TH
Trình bày kiến thức cốt lõi của vận dụng định luật phóng xạ theo số hạt và khối lượng và nêu điều kiện áp dụng.
\loigiai{
Cần nêu được: Số hạt chưa phân rã và khối lượng chất phóng xạ giảm theo $N=N_0 2^{-t/T}$ và $m=m_0 2^{-t/T}$. Khi vận dụng phải xác định đúng dữ kiện, điều kiện áp dụng, hệ đơn vị và kiểm tra tính hợp lí của kết quả. Nhận định “Sau mỗi chu kì bán rã, lượng chất chưa phân rã giảm một lượng cố định như nhau.” sai vì trái với kiến thức trên.
}
\end{bt}

% ===== Câu 51 =====
% ID: L12C4B23-13-TLN
% Mức: TH
\begin{ex}
% ID: L12C4B23-13-TLN
% Mức: TH
Trong bốn phát biểu sau về vận dụng định luật phóng xạ theo số hạt và khối lượng, có bao nhiêu phát biểu đúng? Chỉ ghi một số nguyên. (1) Số hạt chưa phân rã và khối lượng chất phóng xạ giảm theo $N=N_0 2^{-t/T}$ và $m=m_0 2^{-t/T}$. (2) Sau mỗi chu kì bán rã, lượng chất chưa phân rã giảm một lượng cố định như nhau. (3) Cần kiểm tra điều kiện áp dụng trước khi tính toán. (4) Có thể bỏ qua đơn vị của kết quả.
\shortans{2}
\loigiai{
Đối chiếu từng phát biểu với kiến thức cốt lõi: Số hạt chưa phân rã và khối lượng chất phóng xạ giảm theo $N=N_0 2^{-t/T}$ và $m=m_0 2^{-t/T}$. Có 2 phát biểu đúng.
}
\end{ex}

% ===== Câu 52 =====
% ID: L12C4B23-13-TN
% Mức: NB
\dangbt{Viết phương trình phóng xạ và xác định hạt nhân con}
\begin{ex}
% ID: L12C4B23-13-TN
% Mức: NB
Chọn kết luận chính xác về viết phương trình phóng xạ và xác định hạt nhân con?
\choice
{Phóng xạ $\gamma$ làm số khối giảm 1.}
{Có thể bỏ qua điều kiện áp dụng và đơn vị trong mọi trường hợp.}
{\True Phóng xạ $\alpha$ làm $A$ giảm 4, $Z$ giảm 2; $\beta^-$ giữ nguyên $A$ và làm $Z$ tăng 1; $\gamma$ không đổi $A,Z$.}
{Kết quả của bài toán không phụ thuộc dữ kiện đã cho.}
\loigiai{
Phóng xạ $\alpha$ làm $A$ giảm 4, $Z$ giảm 2; $\beta^-$ giữ nguyên $A$ và làm $Z$ tăng 1; $\gamma$ không đổi $A,Z$. Vì vậy phương án $C$ đúng; các phương án còn lại trái với kiến thức hoặc điều kiện áp dụng.
}
\end{ex}

% ===== Câu 53 =====
% ID: L12C4B23-14-DS
% Mức: VD
\dangbt{Vận dụng định luật phóng xạ theo số hạt và khối lượng}
\begin{ex}
% ID: L12C4B23-14-DS
% Mức: VD
Xét vận dụng định luật phóng xạ theo số hạt và khối lượng (tình huống 4). Hãy xác định tính đúng, sai của bốn phát biểu sau.
\choiceTF
{Sau mỗi chu kì bán rã, lượng chất chưa phân rã giảm một lượng cố định như nhau.}
{\True Số hạt chưa phân rã và khối lượng chất phóng xạ giảm theo $N=N_0 2^{-t/T}$ và $m=m_0 2^{-t/T}$.}
{Có thể bỏ qua dữ kiện, đơn vị và ý nghĩa vật lí của kết quả.}
{\True Khi giải cần đọc đúng dữ kiện, dùng hệ đơn vị thống nhất và kiểm tra điều kiện áp dụng.}
\loigiai{
\begin{itemchoice}
\item (Sai) Sau mỗi chu kì bán rã, lượng chất chưa phân rã giảm một lượng cố định như nhau. (Vì): au mỗi chu kì bán rã, lượng chất chưa phân rã giảm một lượng cố định như nhau. b) Đúng: Số hạt chưa phân rã và khối lượng chất phóng xạ giảm theo $N=N_0 2^{-t/T}$ và $m=m_0 2^{-t/T}$. c) Sai: Có thể bỏ qua dữ kiện, đơn vị và ý nghĩa vật lí của kết quả. d) Đúng: Khi giải cần đọc đúng dữ kiện, dùng hệ đơn vị thống nhất và kiểm tra điều kiện áp dụng. Kiến thức cốt lõi: Số hạt chưa phân rã và khối lượng chất phóng xạ giảm theo $N=N_0 2^{-t/T}$ và $m=m_0 2^{-t/T}$
\item (Đúng) Số hạt chưa phân rã và khối lượng chất phóng xạ giảm theo $N=N_0 2^{-t/T}$ và $m=m_0 2^{-t/T}$.
\item (Sai) Có thể bỏ qua dữ kiện, đơn vị và ý nghĩa vật lí của kết quả.
\item (Đúng) Khi giải cần đọc đúng dữ kiện, dùng hệ đơn vị thống nhất và kiểm tra điều kiện áp dụng.
\end{itemchoice}
}
\end{ex}

% ===== Câu 54 =====
% ID: L12C4B23-14-TL
% Mức: TH
\begin{bt}
% ID: L12C4B23-14-TL
% Mức: TH
Giải thích vì sao nhận định sau đúng: “Số hạt chưa phân rã và khối lượng chất phóng xạ giảm theo $N=N_0 2^{-t/T}$ và $m=m_0 2^{-t/T}$.”
\loigiai{
Cần nêu được: Số hạt chưa phân rã và khối lượng chất phóng xạ giảm theo $N=N_0 2^{-t/T}$ và $m=m_0 2^{-t/T}$. Khi vận dụng phải xác định đúng dữ kiện, điều kiện áp dụng, hệ đơn vị và kiểm tra tính hợp lí của kết quả. Nhận định “Sau mỗi chu kì bán rã, lượng chất chưa phân rã giảm một lượng cố định như nhau.” sai vì trái với kiến thức trên.
}
\end{bt}

% ===== Câu 55 =====
% ID: L12C4B23-14-TLN
% Mức: TH
\begin{ex}
% ID: L12C4B23-14-TLN
% Mức: TH
Trong bốn phát biểu sau về vận dụng định luật phóng xạ theo số hạt và khối lượng, có bao nhiêu phát biểu đúng? Chỉ ghi một số nguyên. (1) Sau mỗi chu kì bán rã, lượng chất chưa phân rã giảm một lượng cố định như nhau. (2) Số hạt chưa phân rã và khối lượng chất phóng xạ giảm theo $N=N_0 2^{-t/T}$ và $m=m_0 2^{-t/T}$. (3) Kết quả phải phù hợp ý nghĩa vật lí. (4) Mọi đại lượng đều có cùng bản chất.
\shortans{2}
\loigiai{
Đối chiếu từng phát biểu với kiến thức cốt lõi: Số hạt chưa phân rã và khối lượng chất phóng xạ giảm theo $N=N_0 2^{-t/T}$ và $m=m_0 2^{-t/T}$. Có 2 phát biểu đúng.
}
\end{ex}

% ===== Câu 56 =====
% ID: L12C4B23-14-TN
% Mức: TH
\dangbt{Viết phương trình phóng xạ và xác định hạt nhân con}
\begin{ex}
% ID: L12C4B23-14-TN
% Mức: TH
Nội dung nào mô tả đúng nhất về viết phương trình phóng xạ và xác định hạt nhân con?
\choice
{Phóng xạ $\gamma$ làm số khối giảm 1.}
{Có thể bỏ qua điều kiện áp dụng và đơn vị trong mọi trường hợp.}
{Kết quả của bài toán không phụ thuộc dữ kiện đã cho.}
{\True Phóng xạ $\alpha$ làm $A$ giảm 4, $Z$ giảm 2; $\beta^-$ giữ nguyên $A$ và làm $Z$ tăng 1; $\gamma$ không đổi $A,Z$.}
\loigiai{
Phóng xạ $\alpha$ làm $A$ giảm 4, $Z$ giảm 2; $\beta^-$ giữ nguyên $A$ và làm $Z$ tăng 1; $\gamma$ không đổi $A,Z$. Vì vậy phương án $D$ đúng; các phương án còn lại trái với kiến thức hoặc điều kiện áp dụng.
}
\end{ex}

% ===== Câu 57 =====
% ID: L12C4B23-15-DS
% Mức: VD
\dangbt{Vận dụng định luật phóng xạ theo số hạt và khối lượng}
\begin{ex}
% ID: L12C4B23-15-DS
% Mức: VD
Xét vận dụng định luật phóng xạ theo số hạt và khối lượng (tình huống 5). Hãy xác định tính đúng, sai của bốn phát biểu sau.
\choiceTF
{\True Số hạt chưa phân rã và khối lượng chất phóng xạ giảm theo $N=N_0 2^{-t/T}$ và $m=m_0 2^{-t/T}$.}
{Có thể bỏ qua dữ kiện, đơn vị và ý nghĩa vật lí của kết quả.}
{Sau mỗi chu kì bán rã, lượng chất chưa phân rã giảm một lượng cố định như nhau.}
{Có thể bỏ qua dữ kiện, đơn vị và ý nghĩa vật lí của kết quả.}
\loigiai{
\begin{itemchoice}
\item (Đúng) Số hạt chưa phân rã và khối lượng chất phóng xạ giảm theo $N=N_0 2^{-t/T}$ và $m=m_0 2^{-t/T}$. (Vì): Số hạt chưa phân rã và khối lượng chất phóng xạ giảm theo $N=N_0 2^{-t/T}$ và $m=m_0 2^{-t/T}$. b) Sai: Có thể bỏ qua dữ kiện, đơn vị và ý nghĩa vật lí của kết quả. c) Sai: Sau mỗi chu kì bán rã, lượng chất chưa phân rã giảm một lượng cố định như nhau. d) Sai: Có thể bỏ qua dữ kiện, đơn vị và ý nghĩa vật lí của kết quả. Kiến thức cốt lõi: Số hạt chưa phân rã và khối lượng chất phóng xạ giảm theo $N=N_0 2^{-t/T}$ và $m=m_0 2^{-t/T}$
\item (Sai) Có thể bỏ qua dữ kiện, đơn vị và ý nghĩa vật lí của kết quả.
\item (Sai) Sau mỗi chu kì bán rã, lượng chất chưa phân rã giảm một lượng cố định như nhau.
\item (Sai) Có thể bỏ qua dữ kiện, đơn vị và ý nghĩa vật lí của kết quả.
\end{itemchoice}
}
\end{ex}

% ===== Câu 58 =====
% ID: L12C4B23-15-TL
% Mức: VD
\begin{bt}
% ID: L12C4B23-15-TL
% Mức: VD
Phân tích sai lầm trong nhận định: “Sau mỗi chu kì bán rã, lượng chất chưa phân rã giảm một lượng cố định như nhau.”
\loigiai{
Cần nêu được: Số hạt chưa phân rã và khối lượng chất phóng xạ giảm theo $N=N_0 2^{-t/T}$ và $m=m_0 2^{-t/T}$. Khi vận dụng phải xác định đúng dữ kiện, điều kiện áp dụng, hệ đơn vị và kiểm tra tính hợp lí của kết quả. Nhận định “Sau mỗi chu kì bán rã, lượng chất chưa phân rã giảm một lượng cố định như nhau.” sai vì trái với kiến thức trên.
}
\end{bt}

% ===== Câu 59 =====
% ID: L12C4B23-15-TLN
% Mức: TH
\begin{ex}
% ID: L12C4B23-15-TLN
% Mức: TH
Trong bốn phát biểu sau về vận dụng định luật phóng xạ theo số hạt và khối lượng, có bao nhiêu phát biểu đúng? Chỉ ghi một số nguyên. (1) Cần đổi các đại lượng về hệ đơn vị thống nhất. (2) Số hạt chưa phân rã và khối lượng chất phóng xạ giảm theo $N=N_0 2^{-t/T}$ và $m=m_0 2^{-t/T}$. (3) Sau mỗi chu kì bán rã, lượng chất chưa phân rã giảm một lượng cố định như nhau. (4) Không cần kiểm tra dữ kiện.
\shortans{2}
\loigiai{
Đối chiếu từng phát biểu với kiến thức cốt lõi: Số hạt chưa phân rã và khối lượng chất phóng xạ giảm theo $N=N_0 2^{-t/T}$ và $m=m_0 2^{-t/T}$. Có 2 phát biểu đúng.
}
\end{ex}

% ===== Câu 60 =====
% ID: L12C4B23-15-TN
% Mức: TH
\dangbt{Viết phương trình phóng xạ và xác định hạt nhân con}
\begin{ex}
% ID: L12C4B23-15-TN
% Mức: TH
Khi phân tích, phát biểu nào đúng về viết phương trình phóng xạ và xác định hạt nhân con?
\choice
{\True Phóng xạ $\alpha$ làm $A$ giảm 4, $Z$ giảm 2; $\beta^-$ giữ nguyên $A$ và làm $Z$ tăng 1; $\gamma$ không đổi $A,Z$.}
{Phóng xạ $\gamma$ làm số khối giảm 1.}
{Có thể bỏ qua điều kiện áp dụng và đơn vị trong mọi trường hợp.}
{Kết quả của bài toán không phụ thuộc dữ kiện đã cho.}
\loigiai{
Phóng xạ $\alpha$ làm $A$ giảm 4, $Z$ giảm 2; $\beta^-$ giữ nguyên $A$ và làm $Z$ tăng 1; $\gamma$ không đổi $A,Z$. Vì vậy phương án $A$ đúng; các phương án còn lại trái với kiến thức hoặc điều kiện áp dụng.
}
\end{ex}

% ===== Câu 61 =====
% ID: L12C4B23-16-DS
% Mức: TH
\dangbt{Tính chu kì bán rã, hằng số phóng xạ và tuổi mẫu vật}
\begin{ex}
% ID: L12C4B23-16-DS
% Mức: TH
Xét tính chu kì bán rã, hằng số phóng xạ và tuổi mẫu vật (tình huống 1). Hãy xác định tính đúng, sai của bốn phát biểu sau.
\choiceTF
{\True Hằng số phóng xạ $\lambda=\ln2/T$; từ tỉ số lượng còn lại có thể suy ra thời gian hay tuổi mẫu.}
{Chu kì bán rã càng lớn thì hằng số phóng xạ càng lớn.}
{\True Khi giải cần đọc đúng dữ kiện, dùng hệ đơn vị thống nhất và kiểm tra điều kiện áp dụng.}
{Có thể bỏ qua dữ kiện, đơn vị và ý nghĩa vật lí của kết quả.}
\loigiai{
\begin{itemchoice}
\item (Đúng) Hằng số phóng xạ $\lambda=\ln2/T$; từ tỉ số lượng còn lại có thể suy ra thời gian hay tuổi mẫu. (Vì): Hằng số phóng xạ $\lambda=\ln2/T$; từ tỉ số lượng còn lại có thể suy ra thời gian hay tuổi mẫu. b) Sai: Chu kì bán rã càng lớn thì hằng số phóng xạ càng lớn. c) Đúng: Khi giải cần đọc đúng dữ kiện, dùng hệ đơn vị thống nhất và kiểm tra điều kiện áp dụng. d) Sai: Có thể bỏ qua dữ kiện, đơn vị và ý nghĩa vật lí của kết quả. Kiến thức cốt lõi: Hằng số phóng xạ $\lambda=\ln2/T$; từ tỉ số lượng còn lại có thể suy ra thời gian hay tuổi mẫu
\item (Sai) Chu kì bán rã càng lớn thì hằng số phóng xạ càng lớn.
\item (Đúng) Khi giải cần đọc đúng dữ kiện, dùng hệ đơn vị thống nhất và kiểm tra điều kiện áp dụng.
\item (Sai) Có thể bỏ qua dữ kiện, đơn vị và ý nghĩa vật lí của kết quả.
\end{itemchoice}
}
\end{ex}

% ===== Câu 62 =====
% ID: L12C4B23-16-TL
% Mức: VD
\dangbt{Vận dụng định luật phóng xạ theo số hạt và khối lượng}
\begin{bt}
% ID: L12C4B23-16-TL
% Mức: VD
Đề xuất các bước giải một bài toán thuộc dạng vận dụng định luật phóng xạ theo số hạt và khối lượng.
\loigiai{
Cần nêu được: Số hạt chưa phân rã và khối lượng chất phóng xạ giảm theo $N=N_0 2^{-t/T}$ và $m=m_0 2^{-t/T}$. Khi vận dụng phải xác định đúng dữ kiện, điều kiện áp dụng, hệ đơn vị và kiểm tra tính hợp lí của kết quả. Nhận định “Sau mỗi chu kì bán rã, lượng chất chưa phân rã giảm một lượng cố định như nhau.” sai vì trái với kiến thức trên.
}
\end{bt}

% ===== Câu 63 =====
% ID: L12C4B23-16-TLN
% Mức: VD
\begin{ex}
% ID: L12C4B23-16-TLN
% Mức: VD
Trong bốn phát biểu sau về vận dụng định luật phóng xạ theo số hạt và khối lượng, có bao nhiêu phát biểu đúng? Chỉ ghi một số nguyên. (1) Số hạt chưa phân rã và khối lượng chất phóng xạ giảm theo $N=N_0 2^{-t/T}$ và $m=m_0 2^{-t/T}$. (2) Cần đối chiếu kết quả với điều kiện bài toán. (3) Sau mỗi chu kì bán rã, lượng chất chưa phân rã giảm một lượng cố định như nhau. (4) Mọi trường hợp đều cho cùng một kết quả.
\shortans{2}
\loigiai{
Đối chiếu từng phát biểu với kiến thức cốt lõi: Số hạt chưa phân rã và khối lượng chất phóng xạ giảm theo $N=N_0 2^{-t/T}$ và $m=m_0 2^{-t/T}$. Có 2 phát biểu đúng.
}
\end{ex}

% ===== Câu 64 =====
% ID: L12C4B23-16-TN
% Mức: TH
\dangbt{Viết phương trình phóng xạ và xác định hạt nhân con}
\begin{ex}
% ID: L12C4B23-16-TN
% Mức: TH
Kết luận nào của học sinh là đúng về viết phương trình phóng xạ và xác định hạt nhân con?
\choice
{Phóng xạ $\gamma$ làm số khối giảm 1.}
{\True Phóng xạ $\alpha$ làm $A$ giảm 4, $Z$ giảm 2; $\beta^-$ giữ nguyên $A$ và làm $Z$ tăng 1; $\gamma$ không đổi $A,Z$.}
{Có thể bỏ qua điều kiện áp dụng và đơn vị trong mọi trường hợp.}
{Kết quả của bài toán không phụ thuộc dữ kiện đã cho.}
\loigiai{
Phóng xạ $\alpha$ làm $A$ giảm 4, $Z$ giảm 2; $\beta^-$ giữ nguyên $A$ và làm $Z$ tăng 1; $\gamma$ không đổi $A,Z$. Vì vậy phương án $B$ đúng; các phương án còn lại trái với kiến thức hoặc điều kiện áp dụng.
}
\end{ex}

% ===== Câu 65 =====
% ID: L12C4B23-17-DS
% Mức: TH
\dangbt{Tính chu kì bán rã, hằng số phóng xạ và tuổi mẫu vật}
\begin{ex}
% ID: L12C4B23-17-DS
% Mức: TH
Xét tính chu kì bán rã, hằng số phóng xạ và tuổi mẫu vật (tình huống 2). Hãy xác định tính đúng, sai của bốn phát biểu sau.
\choiceTF
{Chu kì bán rã càng lớn thì hằng số phóng xạ càng lớn.}
{\True Hằng số phóng xạ $\lambda=\ln2/T$; từ tỉ số lượng còn lại có thể suy ra thời gian hay tuổi mẫu.}
{Có thể bỏ qua dữ kiện, đơn vị và ý nghĩa vật lí của kết quả.}
{\True Khi giải cần đọc đúng dữ kiện, dùng hệ đơn vị thống nhất và kiểm tra điều kiện áp dụng.}
\loigiai{
\begin{itemchoice}
\item (Sai) Chu kì bán rã càng lớn thì hằng số phóng xạ càng lớn. (Vì): Chu kì bán rã càng lớn thì hằng số phóng xạ càng lớn. b) Đúng: Hằng số phóng xạ $\lambda=\ln2/T$; từ tỉ số lượng còn lại có thể suy ra thời gian hay tuổi mẫu. c) Sai: Có thể bỏ qua dữ kiện, đơn vị và ý nghĩa vật lí của kết quả. d) Đúng: Khi giải cần đọc đúng dữ kiện, dùng hệ đơn vị thống nhất và kiểm tra điều kiện áp dụng. Kiến thức cốt lõi: Hằng số phóng xạ $\lambda=\ln2/T$; từ tỉ số lượng còn lại có thể suy ra thời gian hay tuổi mẫu
\item (Đúng) Hằng số phóng xạ $\lambda=\ln2/T$; từ tỉ số lượng còn lại có thể suy ra thời gian hay tuổi mẫu.
\item (Sai) Có thể bỏ qua dữ kiện, đơn vị và ý nghĩa vật lí của kết quả.
\item (Đúng) Khi giải cần đọc đúng dữ kiện, dùng hệ đơn vị thống nhất và kiểm tra điều kiện áp dụng.
\end{itemchoice}
}
\end{ex}

% ===== Câu 66 =====
% ID: L12C4B23-17-TL
% Mức: VD
\dangbt{Vận dụng định luật phóng xạ theo số hạt và khối lượng}
\begin{bt}
% ID: L12C4B23-17-TL
% Mức: VD
Nêu một tình huống thực tiễn liên quan đến vận dụng định luật phóng xạ theo số hạt và khối lượng và giải thích bằng kiến thức Vật lí.
\loigiai{
Cần nêu được: Số hạt chưa phân rã và khối lượng chất phóng xạ giảm theo $N=N_0 2^{-t/T}$ và $m=m_0 2^{-t/T}$. Khi vận dụng phải xác định đúng dữ kiện, điều kiện áp dụng, hệ đơn vị và kiểm tra tính hợp lí của kết quả. Nhận định “Sau mỗi chu kì bán rã, lượng chất chưa phân rã giảm một lượng cố định như nhau.” sai vì trái với kiến thức trên.
}
\end{bt}

% ===== Câu 67 =====
% ID: L12C4B23-17-TLN
% Mức: VD
\begin{ex}
% ID: L12C4B23-17-TLN
% Mức: VD
Trong bốn phát biểu sau về vận dụng định luật phóng xạ theo số hạt và khối lượng, có bao nhiêu phát biểu đúng? Chỉ ghi một số nguyên. (1) Sau mỗi chu kì bán rã, lượng chất chưa phân rã giảm một lượng cố định như nhau. (2) Cần đọc đúng dữ kiện và giả thiết. (3) Số hạt chưa phân rã và khối lượng chất phóng xạ giảm theo $N=N_0 2^{-t/T}$ và $m=m_0 2^{-t/T}$. (4) Có thể thay số trước khi chọn công thức.
\shortans{2}
\loigiai{
Đối chiếu từng phát biểu với kiến thức cốt lõi: Số hạt chưa phân rã và khối lượng chất phóng xạ giảm theo $N=N_0 2^{-t/T}$ và $m=m_0 2^{-t/T}$. Có 2 phát biểu đúng.
}
\end{ex}

% ===== Câu 68 =====
% ID: L12C4B23-17-TN
% Mức: TH
\dangbt{Viết phương trình phóng xạ và xác định hạt nhân con}
\begin{ex}
% ID: L12C4B23-17-TN
% Mức: TH
Chọn phát biểu không trái với kiến thức về viết phương trình phóng xạ và xác định hạt nhân con?
\choice
{Phóng xạ $\gamma$ làm số khối giảm 1.}
{Có thể bỏ qua điều kiện áp dụng và đơn vị trong mọi trường hợp.}
{\True Phóng xạ $\alpha$ làm $A$ giảm 4, $Z$ giảm 2; $\beta^-$ giữ nguyên $A$ và làm $Z$ tăng 1; $\gamma$ không đổi $A,Z$.}
{Kết quả của bài toán không phụ thuộc dữ kiện đã cho.}
\loigiai{
Phóng xạ $\alpha$ làm $A$ giảm 4, $Z$ giảm 2; $\beta^-$ giữ nguyên $A$ và làm $Z$ tăng 1; $\gamma$ không đổi $A,Z$. Vì vậy phương án $C$ đúng; các phương án còn lại trái với kiến thức hoặc điều kiện áp dụng.
}
\end{ex}

% ===== Câu 69 =====
% ID: L12C4B23-18-DS
% Mức: VD
\dangbt{Tính chu kì bán rã, hằng số phóng xạ và tuổi mẫu vật}
\begin{ex}
% ID: L12C4B23-18-DS
% Mức: VD
Xét tính chu kì bán rã, hằng số phóng xạ và tuổi mẫu vật (tình huống 3). Hãy xác định tính đúng, sai của bốn phát biểu sau.
\choiceTF
{\True Hằng số phóng xạ $\lambda=\ln2/T$; từ tỉ số lượng còn lại có thể suy ra thời gian hay tuổi mẫu.}
{Chu kì bán rã càng lớn thì hằng số phóng xạ càng lớn.}
{\True Khi giải cần đọc đúng dữ kiện, dùng hệ đơn vị thống nhất và kiểm tra điều kiện áp dụng.}
{Có thể bỏ qua dữ kiện, đơn vị và ý nghĩa vật lí của kết quả.}
\loigiai{
\begin{itemchoice}
\item (Đúng) Hằng số phóng xạ $\lambda=\ln2/T$; từ tỉ số lượng còn lại có thể suy ra thời gian hay tuổi mẫu. (Vì): Hằng số phóng xạ $\lambda=\ln2/T$; từ tỉ số lượng còn lại có thể suy ra thời gian hay tuổi mẫu. b) Sai: Chu kì bán rã càng lớn thì hằng số phóng xạ càng lớn. c) Đúng: Khi giải cần đọc đúng dữ kiện, dùng hệ đơn vị thống nhất và kiểm tra điều kiện áp dụng. d) Sai: Có thể bỏ qua dữ kiện, đơn vị và ý nghĩa vật lí của kết quả. Kiến thức cốt lõi: Hằng số phóng xạ $\lambda=\ln2/T$; từ tỉ số lượng còn lại có thể suy ra thời gian hay tuổi mẫu
\item (Sai) Chu kì bán rã càng lớn thì hằng số phóng xạ càng lớn.
\item (Đúng) Khi giải cần đọc đúng dữ kiện, dùng hệ đơn vị thống nhất và kiểm tra điều kiện áp dụng.
\item (Sai) Có thể bỏ qua dữ kiện, đơn vị và ý nghĩa vật lí của kết quả.
\end{itemchoice}
}
\end{ex}

% ===== Câu 70 =====
% ID: L12C4B23-18-TL
% Mức: VDC
\dangbt{Vận dụng định luật phóng xạ theo số hạt và khối lượng}
\begin{bt}
% ID: L12C4B23-18-TL
% Mức: VDC
So sánh nhận định đúng “Số hạt chưa phân rã và khối lượng chất phóng xạ giảm theo $N=N_0 2^{-t/T}$ và $m=m_0 2^{-t/T}$.” với nhận định sai “Sau mỗi chu kì bán rã, lượng chất chưa phân rã giảm một lượng cố định như nhau.”; chỉ rõ căn cứ Vật lí.
\loigiai{
Cần nêu được: Số hạt chưa phân rã và khối lượng chất phóng xạ giảm theo $N=N_0 2^{-t/T}$ và $m=m_0 2^{-t/T}$. Khi vận dụng phải xác định đúng dữ kiện, điều kiện áp dụng, hệ đơn vị và kiểm tra tính hợp lí của kết quả. Nhận định “Sau mỗi chu kì bán rã, lượng chất chưa phân rã giảm một lượng cố định như nhau.” sai vì trái với kiến thức trên.
}
\end{bt}

% ===== Câu 71 =====
% ID: L12C4B23-18-TLN
% Mức: VDC
\begin{ex}
% ID: L12C4B23-18-TLN
% Mức: VDC
Trong bốn phát biểu sau về vận dụng định luật phóng xạ theo số hạt và khối lượng, có bao nhiêu phát biểu đúng? Chỉ ghi một số nguyên. (1) Kết quả cần có ý nghĩa vật lí. (2) Sau mỗi chu kì bán rã, lượng chất chưa phân rã giảm một lượng cố định như nhau. (3) Phải dùng đúng định luật cho tình huống. (4) Số hạt chưa phân rã và khối lượng chất phóng xạ giảm theo $N=N_0 2^{-t/T}$ và $m=m_0 2^{-t/T}$.
\shortans{3}
\loigiai{
Đối chiếu từng phát biểu với kiến thức cốt lõi: Số hạt chưa phân rã và khối lượng chất phóng xạ giảm theo $N=N_0 2^{-t/T}$ và $m=m_0 2^{-t/T}$. Có 3 phát biểu đúng.
}
\end{ex}

% ===== Câu 72 =====
% ID: L12C4B23-18-TN
% Mức: VD
\dangbt{Viết phương trình phóng xạ và xác định hạt nhân con}
\begin{ex}
% ID: L12C4B23-18-TN
% Mức: VD
Cơ sở Vật lí đúng để giải bài là gì về viết phương trình phóng xạ và xác định hạt nhân con?
\choice
{Phóng xạ $\gamma$ làm số khối giảm 1.}
{Có thể bỏ qua điều kiện áp dụng và đơn vị trong mọi trường hợp.}
{Kết quả của bài toán không phụ thuộc dữ kiện đã cho.}
{\True Phóng xạ $\alpha$ làm $A$ giảm 4, $Z$ giảm 2; $\beta^-$ giữ nguyên $A$ và làm $Z$ tăng 1; $\gamma$ không đổi $A,Z$.}
\loigiai{
Phóng xạ $\alpha$ làm $A$ giảm 4, $Z$ giảm 2; $\beta^-$ giữ nguyên $A$ và làm $Z$ tăng 1; $\gamma$ không đổi $A,Z$. Vì vậy phương án $D$ đúng; các phương án còn lại trái với kiến thức hoặc điều kiện áp dụng.
}
\end{ex}

% ===== Câu 73 =====
% ID: L12C4B23-19-DS
% Mức: VD
\dangbt{Tính chu kì bán rã, hằng số phóng xạ và tuổi mẫu vật}
\begin{ex}
% ID: L12C4B23-19-DS
% Mức: VD
Xét tính chu kì bán rã, hằng số phóng xạ và tuổi mẫu vật (tình huống 4). Hãy xác định tính đúng, sai của bốn phát biểu sau.
\choiceTF
{Chu kì bán rã càng lớn thì hằng số phóng xạ càng lớn.}
{\True Hằng số phóng xạ $\lambda=\ln2/T$; từ tỉ số lượng còn lại có thể suy ra thời gian hay tuổi mẫu.}
{Có thể bỏ qua dữ kiện, đơn vị và ý nghĩa vật lí của kết quả.}
{\True Khi giải cần đọc đúng dữ kiện, dùng hệ đơn vị thống nhất và kiểm tra điều kiện áp dụng.}
\loigiai{
\begin{itemchoice}
\item (Sai) Chu kì bán rã càng lớn thì hằng số phóng xạ càng lớn. (Vì): Chu kì bán rã càng lớn thì hằng số phóng xạ càng lớn. b) Đúng: Hằng số phóng xạ $\lambda=\ln2/T$; từ tỉ số lượng còn lại có thể suy ra thời gian hay tuổi mẫu. c) Sai: Có thể bỏ qua dữ kiện, đơn vị và ý nghĩa vật lí của kết quả. d) Đúng: Khi giải cần đọc đúng dữ kiện, dùng hệ đơn vị thống nhất và kiểm tra điều kiện áp dụng. Kiến thức cốt lõi: Hằng số phóng xạ $\lambda=\ln2/T$; từ tỉ số lượng còn lại có thể suy ra thời gian hay tuổi mẫu
\item (Đúng) Hằng số phóng xạ $\lambda=\ln2/T$; từ tỉ số lượng còn lại có thể suy ra thời gian hay tuổi mẫu.
\item (Sai) Có thể bỏ qua dữ kiện, đơn vị và ý nghĩa vật lí của kết quả.
\item (Đúng) Khi giải cần đọc đúng dữ kiện, dùng hệ đơn vị thống nhất và kiểm tra điều kiện áp dụng.
\end{itemchoice}
}
\end{ex}

% ===== Câu 74 =====
% ID: L12C4B23-19-TL
% Mức: TH
\begin{bt}
% ID: L12C4B23-19-TL
% Mức: TH
Trình bày kiến thức cốt lõi của tính chu kì bán rã, hằng số phóng xạ và tuổi mẫu vật và nêu điều kiện áp dụng.
\loigiai{
Cần nêu được: Hằng số phóng xạ $\lambda=\ln2/T$; từ tỉ số lượng còn lại có thể suy ra thời gian hay tuổi mẫu. Khi vận dụng phải xác định đúng dữ kiện, điều kiện áp dụng, hệ đơn vị và kiểm tra tính hợp lí của kết quả. Nhận định “Chu kì bán rã càng lớn thì hằng số phóng xạ càng lớn.” sai vì trái với kiến thức trên.
}
\end{bt}

% ===== Câu 75 =====
% ID: L12C4B23-19-TLN
% Mức: TH
\begin{ex}
% ID: L12C4B23-19-TLN
% Mức: TH
Trong bốn phát biểu sau về tính chu kì bán rã, hằng số phóng xạ và tuổi mẫu vật, có bao nhiêu phát biểu đúng? Chỉ ghi một số nguyên. (1) Hằng số phóng xạ $\lambda=\ln2/T$; từ tỉ số lượng còn lại có thể suy ra thời gian hay tuổi mẫu. (2) Chu kì bán rã càng lớn thì hằng số phóng xạ càng lớn. (3) Cần kiểm tra điều kiện áp dụng trước khi tính toán. (4) Có thể bỏ qua đơn vị của kết quả.
\shortans{2}
\loigiai{
Đối chiếu từng phát biểu với kiến thức cốt lõi: Hằng số phóng xạ $\lambda=\ln2/T$; từ tỉ số lượng còn lại có thể suy ra thời gian hay tuổi mẫu. Có 2 phát biểu đúng.
}
\end{ex}

% ===== Câu 76 =====
% ID: L12C4B23-19-TN
% Mức: VD
\dangbt{Viết phương trình phóng xạ và xác định hạt nhân con}
\begin{ex}
% ID: L12C4B23-19-TN
% Mức: VD
Trong một bài thực hành, nhận xét nào đúng về viết phương trình phóng xạ và xác định hạt nhân con?
\choice
{\True Phóng xạ $\alpha$ làm $A$ giảm 4, $Z$ giảm 2; $\beta^-$ giữ nguyên $A$ và làm $Z$ tăng 1; $\gamma$ không đổi $A,Z$.}
{Phóng xạ $\gamma$ làm số khối giảm 1.}
{Có thể bỏ qua điều kiện áp dụng và đơn vị trong mọi trường hợp.}
{Kết quả của bài toán không phụ thuộc dữ kiện đã cho.}
\loigiai{
Phóng xạ $\alpha$ làm $A$ giảm 4, $Z$ giảm 2; $\beta^-$ giữ nguyên $A$ và làm $Z$ tăng 1; $\gamma$ không đổi $A,Z$. Vì vậy phương án $A$ đúng; các phương án còn lại trái với kiến thức hoặc điều kiện áp dụng.
}
\end{ex}

% ===== Câu 77 =====
% ID: L12C4B23-20-DS
% Mức: VD
\dangbt{Tính chu kì bán rã, hằng số phóng xạ và tuổi mẫu vật}
\begin{ex}
% ID: L12C4B23-20-DS
% Mức: VD
Xét tính chu kì bán rã, hằng số phóng xạ và tuổi mẫu vật (tình huống 5). Hãy xác định tính đúng, sai của bốn phát biểu sau.
\choiceTF
{\True Hằng số phóng xạ $\lambda=\ln2/T$; từ tỉ số lượng còn lại có thể suy ra thời gian hay tuổi mẫu.}
{Có thể bỏ qua dữ kiện, đơn vị và ý nghĩa vật lí của kết quả.}
{Chu kì bán rã càng lớn thì hằng số phóng xạ càng lớn.}
{Có thể bỏ qua dữ kiện, đơn vị và ý nghĩa vật lí của kết quả.}
\loigiai{
\begin{itemchoice}
\item (Đúng) Hằng số phóng xạ $\lambda=\ln2/T$; từ tỉ số lượng còn lại có thể suy ra thời gian hay tuổi mẫu. (Vì): Hằng số phóng xạ $\lambda=\ln2/T$; từ tỉ số lượng còn lại có thể suy ra thời gian hay tuổi mẫu. b) Sai: Có thể bỏ qua dữ kiện, đơn vị và ý nghĩa vật lí của kết quả. c) Sai: Chu kì bán rã càng lớn thì hằng số phóng xạ càng lớn. d) Sai: Có thể bỏ qua dữ kiện, đơn vị và ý nghĩa vật lí của kết quả. Kiến thức cốt lõi: Hằng số phóng xạ $\lambda=\ln2/T$; từ tỉ số lượng còn lại có thể suy ra thời gian hay tuổi mẫu
\item (Sai) Có thể bỏ qua dữ kiện, đơn vị và ý nghĩa vật lí của kết quả.
\item (Sai) Chu kì bán rã càng lớn thì hằng số phóng xạ càng lớn.
\item (Sai) Có thể bỏ qua dữ kiện, đơn vị và ý nghĩa vật lí của kết quả.
\end{itemchoice}
}
\end{ex}

% ===== Câu 78 =====
% ID: L12C4B23-20-TL
% Mức: TH
\begin{bt}
% ID: L12C4B23-20-TL
% Mức: TH
Giải thích vì sao nhận định sau đúng: “Hằng số phóng xạ $\lambda=\ln2/T$; từ tỉ số lượng còn lại có thể suy ra thời gian hay tuổi mẫu.”
\loigiai{
Cần nêu được: Hằng số phóng xạ $\lambda=\ln2/T$; từ tỉ số lượng còn lại có thể suy ra thời gian hay tuổi mẫu. Khi vận dụng phải xác định đúng dữ kiện, điều kiện áp dụng, hệ đơn vị và kiểm tra tính hợp lí của kết quả. Nhận định “Chu kì bán rã càng lớn thì hằng số phóng xạ càng lớn.” sai vì trái với kiến thức trên.
}
\end{bt}

% ===== Câu 79 =====
% ID: L12C4B23-20-TLN
% Mức: TH
\begin{ex}
% ID: L12C4B23-20-TLN
% Mức: TH
Trong bốn phát biểu sau về tính chu kì bán rã, hằng số phóng xạ và tuổi mẫu vật, có bao nhiêu phát biểu đúng? Chỉ ghi một số nguyên. (1) Chu kì bán rã càng lớn thì hằng số phóng xạ càng lớn. (2) Hằng số phóng xạ $\lambda=\ln2/T$; từ tỉ số lượng còn lại có thể suy ra thời gian hay tuổi mẫu. (3) Kết quả phải phù hợp ý nghĩa vật lí. (4) Mọi đại lượng đều có cùng bản chất.
\shortans{2}
\loigiai{
Đối chiếu từng phát biểu với kiến thức cốt lõi: Hằng số phóng xạ $\lambda=\ln2/T$; từ tỉ số lượng còn lại có thể suy ra thời gian hay tuổi mẫu. Có 2 phát biểu đúng.
}
\end{ex}

% ===== Câu 80 =====
% ID: L12C4B23-20-TN
% Mức: VD
\dangbt{Viết phương trình phóng xạ và xác định hạt nhân con}
\begin{ex}
% ID: L12C4B23-20-TN
% Mức: VD
Kết luận đúng khi vận dụng là về viết phương trình phóng xạ và xác định hạt nhân con?
\choice
{Phóng xạ $\gamma$ làm số khối giảm 1.}
{\True Phóng xạ $\alpha$ làm $A$ giảm 4, $Z$ giảm 2; $\beta^-$ giữ nguyên $A$ và làm $Z$ tăng 1; $\gamma$ không đổi $A,Z$.}
{Có thể bỏ qua điều kiện áp dụng và đơn vị trong mọi trường hợp.}
{Kết quả của bài toán không phụ thuộc dữ kiện đã cho.}
\loigiai{
Phóng xạ $\alpha$ làm $A$ giảm 4, $Z$ giảm 2; $\beta^-$ giữ nguyên $A$ và làm $Z$ tăng 1; $\gamma$ không đổi $A,Z$. Vì vậy phương án $B$ đúng; các phương án còn lại trái với kiến thức hoặc điều kiện áp dụng.
}
\end{ex}

% ===== Câu 81 =====
% ID: L12C4B23-21-DS
% Mức: TH
\dangbt{Tính độ phóng xạ và khai thác đồ thị phóng xạ}
\begin{ex}
% ID: L12C4B23-21-DS
% Mức: TH
Xét tính độ phóng xạ và khai thác đồ thị phóng xạ (tình huống 1). Hãy xác định tính đúng, sai của bốn phát biểu sau.
\choiceTF
{\True Độ phóng xạ $H=\lambda N$ và cũng giảm theo quy luật mũ với cùng chu kì bán rã.}
{Độ phóng xạ không phụ thuộc số hạt chưa phân rã.}
{\True Khi giải cần đọc đúng dữ kiện, dùng hệ đơn vị thống nhất và kiểm tra điều kiện áp dụng.}
{Có thể bỏ qua dữ kiện, đơn vị và ý nghĩa vật lí của kết quả.}
\loigiai{
\begin{itemchoice}
\item (Đúng) Độ phóng xạ $H=\lambda N$ và cũng giảm theo quy luật mũ với cùng chu kì bán rã. (Vì): ộ phóng xạ $H=\lambda N$ và cũng giảm theo quy luật mũ với cùng chu kì bán rã. b) Sai: Độ phóng xạ không phụ thuộc số hạt chưa phân rã. c) Đúng: Khi giải cần đọc đúng dữ kiện, dùng hệ đơn vị thống nhất và kiểm tra điều kiện áp dụng. d) Sai: Có thể bỏ qua dữ kiện, đơn vị và ý nghĩa vật lí của kết quả. Kiến thức cốt lõi: Độ phóng xạ $H=\lambda N$ và cũng giảm theo quy luật mũ với cùng chu kì bán rã
\item (Sai) Độ phóng xạ không phụ thuộc số hạt chưa phân rã.
\item (Đúng) Khi giải cần đọc đúng dữ kiện, dùng hệ đơn vị thống nhất và kiểm tra điều kiện áp dụng.
\item (Sai) Có thể bỏ qua dữ kiện, đơn vị và ý nghĩa vật lí của kết quả.
\end{itemchoice}
}
\end{ex}

% ===== Câu 82 =====
% ID: L12C4B23-21-TL
% Mức: VD
\dangbt{Tính chu kì bán rã, hằng số phóng xạ và tuổi mẫu vật}
\begin{bt}
% ID: L12C4B23-21-TL
% Mức: VD
Phân tích sai lầm trong nhận định: “Chu kì bán rã càng lớn thì hằng số phóng xạ càng lớn.”
\loigiai{
Cần nêu được: Hằng số phóng xạ $\lambda=\ln2/T$; từ tỉ số lượng còn lại có thể suy ra thời gian hay tuổi mẫu. Khi vận dụng phải xác định đúng dữ kiện, điều kiện áp dụng, hệ đơn vị và kiểm tra tính hợp lí của kết quả. Nhận định “Chu kì bán rã càng lớn thì hằng số phóng xạ càng lớn.” sai vì trái với kiến thức trên.
}
\end{bt}

% ===== Câu 83 =====
% ID: L12C4B23-21-TLN
% Mức: TH
\begin{ex}
% ID: L12C4B23-21-TLN
% Mức: TH
Trong bốn phát biểu sau về tính chu kì bán rã, hằng số phóng xạ và tuổi mẫu vật, có bao nhiêu phát biểu đúng? Chỉ ghi một số nguyên. (1) Cần đổi các đại lượng về hệ đơn vị thống nhất. (2) Hằng số phóng xạ $\lambda=\ln2/T$; từ tỉ số lượng còn lại có thể suy ra thời gian hay tuổi mẫu. (3) Chu kì bán rã càng lớn thì hằng số phóng xạ càng lớn. (4) Không cần kiểm tra dữ kiện.
\shortans{2}
\loigiai{
Đối chiếu từng phát biểu với kiến thức cốt lõi: Hằng số phóng xạ $\lambda=\ln2/T$; từ tỉ số lượng còn lại có thể suy ra thời gian hay tuổi mẫu. Có 2 phát biểu đúng.
}
\end{ex}

% ===== Câu 84 =====
% ID: L12C4B23-21-TN
% Mức: NB
\dangbt{Vận dụng định luật phóng xạ theo số hạt và khối lượng}
\begin{ex}
% ID: L12C4B23-21-TN
% Mức: NB
Phát biểu nào sau đây đúng về vận dụng định luật phóng xạ theo số hạt và khối lượng?
\choice
{\True Số hạt chưa phân rã và khối lượng chất phóng xạ giảm theo $N=N_0 2^{-t/T}$ và $m=m_0 2^{-t/T}$.}
{Sau mỗi chu kì bán rã, lượng chất chưa phân rã giảm một lượng cố định như nhau.}
{Có thể bỏ qua điều kiện áp dụng và đơn vị trong mọi trường hợp.}
{Kết quả của bài toán không phụ thuộc dữ kiện đã cho.}
\loigiai{
Số hạt chưa phân rã và khối lượng chất phóng xạ giảm theo $N=N_0 2^{-t/T}$ và $m=m_0 2^{-t/T}$. Vì vậy phương án $A$ đúng; các phương án còn lại trái với kiến thức hoặc điều kiện áp dụng.
}
\end{ex}

% ===== Câu 85 =====
% ID: L12C4B23-22-DS
% Mức: TH
\dangbt{Tính độ phóng xạ và khai thác đồ thị phóng xạ}
\begin{ex}
% ID: L12C4B23-22-DS
% Mức: TH
Xét tính độ phóng xạ và khai thác đồ thị phóng xạ (tình huống 2). Hãy xác định tính đúng, sai của bốn phát biểu sau.
\choiceTF
{Độ phóng xạ không phụ thuộc số hạt chưa phân rã.}
{\True Độ phóng xạ $H=\lambda N$ và cũng giảm theo quy luật mũ với cùng chu kì bán rã.}
{Có thể bỏ qua dữ kiện, đơn vị và ý nghĩa vật lí của kết quả.}
{\True Khi giải cần đọc đúng dữ kiện, dùng hệ đơn vị thống nhất và kiểm tra điều kiện áp dụng.}
\loigiai{
\begin{itemchoice}
\item (Sai) Độ phóng xạ không phụ thuộc số hạt chưa phân rã. (Vì): Độ phóng xạ không phụ thuộc số hạt chưa phân rã. b) Đúng: Độ phóng xạ $H=\lambda N$ và cũng giảm theo quy luật mũ với cùng chu kì bán rã. c) Sai: Có thể bỏ qua dữ kiện, đơn vị và ý nghĩa vật lí của kết quả. d) Đúng: Khi giải cần đọc đúng dữ kiện, dùng hệ đơn vị thống nhất và kiểm tra điều kiện áp dụng. Kiến thức cốt lõi: Độ phóng xạ $H=\lambda N$ và cũng giảm theo quy luật mũ với cùng chu kì bán rã
\item (Đúng) Độ phóng xạ $H=\lambda N$ và cũng giảm theo quy luật mũ với cùng chu kì bán rã.
\item (Sai) Có thể bỏ qua dữ kiện, đơn vị và ý nghĩa vật lí của kết quả.
\item (Đúng) Khi giải cần đọc đúng dữ kiện, dùng hệ đơn vị thống nhất và kiểm tra điều kiện áp dụng.
\end{itemchoice}
}
\end{ex}

% ===== Câu 86 =====
% ID: L12C4B23-22-TL
% Mức: VD
\dangbt{Tính chu kì bán rã, hằng số phóng xạ và tuổi mẫu vật}
\begin{bt}
% ID: L12C4B23-22-TL
% Mức: VD
Đề xuất các bước giải một bài toán thuộc dạng tính chu kì bán rã, hằng số phóng xạ và tuổi mẫu vật.
\loigiai{
Cần nêu được: Hằng số phóng xạ $\lambda=\ln2/T$; từ tỉ số lượng còn lại có thể suy ra thời gian hay tuổi mẫu. Khi vận dụng phải xác định đúng dữ kiện, điều kiện áp dụng, hệ đơn vị và kiểm tra tính hợp lí của kết quả. Nhận định “Chu kì bán rã càng lớn thì hằng số phóng xạ càng lớn.” sai vì trái với kiến thức trên.
}
\end{bt}

% ===== Câu 87 =====
% ID: L12C4B23-22-TLN
% Mức: VD
\begin{ex}
% ID: L12C4B23-22-TLN
% Mức: VD
Trong bốn phát biểu sau về tính chu kì bán rã, hằng số phóng xạ và tuổi mẫu vật, có bao nhiêu phát biểu đúng? Chỉ ghi một số nguyên. (1) Hằng số phóng xạ $\lambda=\ln2/T$; từ tỉ số lượng còn lại có thể suy ra thời gian hay tuổi mẫu. (2) Cần đối chiếu kết quả với điều kiện bài toán. (3) Chu kì bán rã càng lớn thì hằng số phóng xạ càng lớn. (4) Mọi trường hợp đều cho cùng một kết quả.
\shortans{2}
\loigiai{
Đối chiếu từng phát biểu với kiến thức cốt lõi: Hằng số phóng xạ $\lambda=\ln2/T$; từ tỉ số lượng còn lại có thể suy ra thời gian hay tuổi mẫu. Có 2 phát biểu đúng.
}
\end{ex}

% ===== Câu 88 =====
% ID: L12C4B23-22-TN
% Mức: NB
\dangbt{Vận dụng định luật phóng xạ theo số hạt và khối lượng}
\begin{ex}
% ID: L12C4B23-22-TN
% Mức: NB
Nhận định nào phù hợp với kiến thức Vật lí về vận dụng định luật phóng xạ theo số hạt và khối lượng?
\choice
{Sau mỗi chu kì bán rã, lượng chất chưa phân rã giảm một lượng cố định như nhau.}
{\True Số hạt chưa phân rã và khối lượng chất phóng xạ giảm theo $N=N_0 2^{-t/T}$ và $m=m_0 2^{-t/T}$.}
{Có thể bỏ qua điều kiện áp dụng và đơn vị trong mọi trường hợp.}
{Kết quả của bài toán không phụ thuộc dữ kiện đã cho.}
\loigiai{
Số hạt chưa phân rã và khối lượng chất phóng xạ giảm theo $N=N_0 2^{-t/T}$ và $m=m_0 2^{-t/T}$. Vì vậy phương án $B$ đúng; các phương án còn lại trái với kiến thức hoặc điều kiện áp dụng.
}
\end{ex}

% ===== Câu 89 =====
% ID: L12C4B23-23-DS
% Mức: VD
\dangbt{Tính độ phóng xạ và khai thác đồ thị phóng xạ}
\begin{ex}
% ID: L12C4B23-23-DS
% Mức: VD
Xét tính độ phóng xạ và khai thác đồ thị phóng xạ (tình huống 3). Hãy xác định tính đúng, sai của bốn phát biểu sau.
\choiceTF
{\True Độ phóng xạ $H=\lambda N$ và cũng giảm theo quy luật mũ với cùng chu kì bán rã.}
{Độ phóng xạ không phụ thuộc số hạt chưa phân rã.}
{\True Khi giải cần đọc đúng dữ kiện, dùng hệ đơn vị thống nhất và kiểm tra điều kiện áp dụng.}
{Có thể bỏ qua dữ kiện, đơn vị và ý nghĩa vật lí của kết quả.}
\loigiai{
\begin{itemchoice}
\item (Đúng) Độ phóng xạ $H=\lambda N$ và cũng giảm theo quy luật mũ với cùng chu kì bán rã. (Vì): ộ phóng xạ $H=\lambda N$ và cũng giảm theo quy luật mũ với cùng chu kì bán rã. b) Sai: Độ phóng xạ không phụ thuộc số hạt chưa phân rã. c) Đúng: Khi giải cần đọc đúng dữ kiện, dùng hệ đơn vị thống nhất và kiểm tra điều kiện áp dụng. d) Sai: Có thể bỏ qua dữ kiện, đơn vị và ý nghĩa vật lí của kết quả. Kiến thức cốt lõi: Độ phóng xạ $H=\lambda N$ và cũng giảm theo quy luật mũ với cùng chu kì bán rã
\item (Sai) Độ phóng xạ không phụ thuộc số hạt chưa phân rã.
\item (Đúng) Khi giải cần đọc đúng dữ kiện, dùng hệ đơn vị thống nhất và kiểm tra điều kiện áp dụng.
\item (Sai) Có thể bỏ qua dữ kiện, đơn vị và ý nghĩa vật lí của kết quả.
\end{itemchoice}
}
\end{ex}

% ===== Câu 90 =====
% ID: L12C4B23-23-TL
% Mức: VD
\dangbt{Tính chu kì bán rã, hằng số phóng xạ và tuổi mẫu vật}
\begin{bt}
% ID: L12C4B23-23-TL
% Mức: VD
Nêu một tình huống thực tiễn liên quan đến tính chu kì bán rã, hằng số phóng xạ và tuổi mẫu vật và giải thích bằng kiến thức Vật lí.
\loigiai{
Cần nêu được: Hằng số phóng xạ $\lambda=\ln2/T$; từ tỉ số lượng còn lại có thể suy ra thời gian hay tuổi mẫu. Khi vận dụng phải xác định đúng dữ kiện, điều kiện áp dụng, hệ đơn vị và kiểm tra tính hợp lí của kết quả. Nhận định “Chu kì bán rã càng lớn thì hằng số phóng xạ càng lớn.” sai vì trái với kiến thức trên.
}
\end{bt}

% ===== Câu 91 =====
% ID: L12C4B23-23-TLN
% Mức: VD
\begin{ex}
% ID: L12C4B23-23-TLN
% Mức: VD
Trong bốn phát biểu sau về tính chu kì bán rã, hằng số phóng xạ và tuổi mẫu vật, có bao nhiêu phát biểu đúng? Chỉ ghi một số nguyên. (1) Chu kì bán rã càng lớn thì hằng số phóng xạ càng lớn. (2) Cần đọc đúng dữ kiện và giả thiết. (3) Hằng số phóng xạ $\lambda=\ln2/T$; từ tỉ số lượng còn lại có thể suy ra thời gian hay tuổi mẫu. (4) Có thể thay số trước khi chọn công thức.
\shortans{2}
\loigiai{
Đối chiếu từng phát biểu với kiến thức cốt lõi: Hằng số phóng xạ $\lambda=\ln2/T$; từ tỉ số lượng còn lại có thể suy ra thời gian hay tuổi mẫu. Có 2 phát biểu đúng.
}
\end{ex}

% ===== Câu 92 =====
% ID: L12C4B23-23-TN
% Mức: NB
\dangbt{Vận dụng định luật phóng xạ theo số hạt và khối lượng}
\begin{ex}
% ID: L12C4B23-23-TN
% Mức: NB
Chọn kết luận chính xác về vận dụng định luật phóng xạ theo số hạt và khối lượng?
\choice
{Sau mỗi chu kì bán rã, lượng chất chưa phân rã giảm một lượng cố định như nhau.}
{Có thể bỏ qua điều kiện áp dụng và đơn vị trong mọi trường hợp.}
{\True Số hạt chưa phân rã và khối lượng chất phóng xạ giảm theo $N=N_0 2^{-t/T}$ và $m=m_0 2^{-t/T}$.}
{Kết quả của bài toán không phụ thuộc dữ kiện đã cho.}
\loigiai{
Số hạt chưa phân rã và khối lượng chất phóng xạ giảm theo $N=N_0 2^{-t/T}$ và $m=m_0 2^{-t/T}$. Vì vậy phương án $C$ đúng; các phương án còn lại trái với kiến thức hoặc điều kiện áp dụng.
}
\end{ex}

% ===== Câu 93 =====
% ID: L12C4B23-24-DS
% Mức: VD
\dangbt{Tính độ phóng xạ và khai thác đồ thị phóng xạ}
\begin{ex}
% ID: L12C4B23-24-DS
% Mức: VD
Xét tính độ phóng xạ và khai thác đồ thị phóng xạ (tình huống 4). Hãy xác định tính đúng, sai của bốn phát biểu sau.
\choiceTF
{Độ phóng xạ không phụ thuộc số hạt chưa phân rã.}
{\True Độ phóng xạ $H=\lambda N$ và cũng giảm theo quy luật mũ với cùng chu kì bán rã.}
{Có thể bỏ qua dữ kiện, đơn vị và ý nghĩa vật lí của kết quả.}
{\True Khi giải cần đọc đúng dữ kiện, dùng hệ đơn vị thống nhất và kiểm tra điều kiện áp dụng.}
\loigiai{
\begin{itemchoice}
\item (Sai) Độ phóng xạ không phụ thuộc số hạt chưa phân rã. (Vì): Độ phóng xạ không phụ thuộc số hạt chưa phân rã. b) Đúng: Độ phóng xạ $H=\lambda N$ và cũng giảm theo quy luật mũ với cùng chu kì bán rã. c) Sai: Có thể bỏ qua dữ kiện, đơn vị và ý nghĩa vật lí của kết quả. d) Đúng: Khi giải cần đọc đúng dữ kiện, dùng hệ đơn vị thống nhất và kiểm tra điều kiện áp dụng. Kiến thức cốt lõi: Độ phóng xạ $H=\lambda N$ và cũng giảm theo quy luật mũ với cùng chu kì bán rã
\item (Đúng) Độ phóng xạ $H=\lambda N$ và cũng giảm theo quy luật mũ với cùng chu kì bán rã.
\item (Sai) Có thể bỏ qua dữ kiện, đơn vị và ý nghĩa vật lí của kết quả.
\item (Đúng) Khi giải cần đọc đúng dữ kiện, dùng hệ đơn vị thống nhất và kiểm tra điều kiện áp dụng.
\end{itemchoice}
}
\end{ex}

% ===== Câu 94 =====
% ID: L12C4B23-24-TL
% Mức: VDC
\dangbt{Tính chu kì bán rã, hằng số phóng xạ và tuổi mẫu vật}
\begin{bt}
% ID: L12C4B23-24-TL
% Mức: VDC
So sánh nhận định đúng “Hằng số phóng xạ $\lambda=\ln2/T$; từ tỉ số lượng còn lại có thể suy ra thời gian hay tuổi mẫu.” với nhận định sai “Chu kì bán rã càng lớn thì hằng số phóng xạ càng lớn.”; chỉ rõ căn cứ Vật lí.
\loigiai{
Cần nêu được: Hằng số phóng xạ $\lambda=\ln2/T$; từ tỉ số lượng còn lại có thể suy ra thời gian hay tuổi mẫu. Khi vận dụng phải xác định đúng dữ kiện, điều kiện áp dụng, hệ đơn vị và kiểm tra tính hợp lí của kết quả. Nhận định “Chu kì bán rã càng lớn thì hằng số phóng xạ càng lớn.” sai vì trái với kiến thức trên.
}
\end{bt}

% ===== Câu 95 =====
% ID: L12C4B23-24-TLN
% Mức: VDC
\begin{ex}
% ID: L12C4B23-24-TLN
% Mức: VDC
Trong bốn phát biểu sau về tính chu kì bán rã, hằng số phóng xạ và tuổi mẫu vật, có bao nhiêu phát biểu đúng? Chỉ ghi một số nguyên. (1) Kết quả cần có ý nghĩa vật lí. (2) Chu kì bán rã càng lớn thì hằng số phóng xạ càng lớn. (3) Phải dùng đúng định luật cho tình huống. (4) Hằng số phóng xạ $\lambda=\ln2/T$; từ tỉ số lượng còn lại có thể suy ra thời gian hay tuổi mẫu.
\shortans{3}
\loigiai{
Đối chiếu từng phát biểu với kiến thức cốt lõi: Hằng số phóng xạ $\lambda=\ln2/T$; từ tỉ số lượng còn lại có thể suy ra thời gian hay tuổi mẫu. Có 3 phát biểu đúng.
}
\end{ex}

% ===== Câu 96 =====
% ID: L12C4B23-24-TN
% Mức: TH
\dangbt{Vận dụng định luật phóng xạ theo số hạt và khối lượng}
\begin{ex}
% ID: L12C4B23-24-TN
% Mức: TH
Nội dung nào mô tả đúng nhất về vận dụng định luật phóng xạ theo số hạt và khối lượng?
\choice
{Sau mỗi chu kì bán rã, lượng chất chưa phân rã giảm một lượng cố định như nhau.}
{Có thể bỏ qua điều kiện áp dụng và đơn vị trong mọi trường hợp.}
{Kết quả của bài toán không phụ thuộc dữ kiện đã cho.}
{\True Số hạt chưa phân rã và khối lượng chất phóng xạ giảm theo $N=N_0 2^{-t/T}$ và $m=m_0 2^{-t/T}$.}
\loigiai{
Số hạt chưa phân rã và khối lượng chất phóng xạ giảm theo $N=N_0 2^{-t/T}$ và $m=m_0 2^{-t/T}$. Vì vậy phương án $D$ đúng; các phương án còn lại trái với kiến thức hoặc điều kiện áp dụng.
}
\end{ex}

% ===== Câu 97 =====
% ID: L12C4B23-25-DS
% Mức: VD
\dangbt{Tính độ phóng xạ và khai thác đồ thị phóng xạ}
\begin{ex}
% ID: L12C4B23-25-DS
% Mức: VD
Xét tính độ phóng xạ và khai thác đồ thị phóng xạ (tình huống 5). Hãy xác định tính đúng, sai của bốn phát biểu sau.
\choiceTF
{\True Độ phóng xạ $H=\lambda N$ và cũng giảm theo quy luật mũ với cùng chu kì bán rã.}
{Có thể bỏ qua dữ kiện, đơn vị và ý nghĩa vật lí của kết quả.}
{Độ phóng xạ không phụ thuộc số hạt chưa phân rã.}
{Có thể bỏ qua dữ kiện, đơn vị và ý nghĩa vật lí của kết quả.}
\loigiai{
\begin{itemchoice}
\item (Đúng) Độ phóng xạ $H=\lambda N$ và cũng giảm theo quy luật mũ với cùng chu kì bán rã. (Vì): ộ phóng xạ $H=\lambda N$ và cũng giảm theo quy luật mũ với cùng chu kì bán rã. b) Sai: Có thể bỏ qua dữ kiện, đơn vị và ý nghĩa vật lí của kết quả. c) Sai: Độ phóng xạ không phụ thuộc số hạt chưa phân rã. d) Sai: Có thể bỏ qua dữ kiện, đơn vị và ý nghĩa vật lí của kết quả. Kiến thức cốt lõi: Độ phóng xạ $H=\lambda N$ và cũng giảm theo quy luật mũ với cùng chu kì bán rã
\item (Sai) Có thể bỏ qua dữ kiện, đơn vị và ý nghĩa vật lí của kết quả.
\item (Sai) Độ phóng xạ không phụ thuộc số hạt chưa phân rã.
\item (Sai) Có thể bỏ qua dữ kiện, đơn vị và ý nghĩa vật lí của kết quả.
\end{itemchoice}
}
\end{ex}

% ===== Câu 98 =====
% ID: L12C4B23-25-TL
% Mức: TH
\begin{bt}
% ID: L12C4B23-25-TL
% Mức: TH
Trình bày kiến thức cốt lõi của tính độ phóng xạ và khai thác đồ thị phóng xạ và nêu điều kiện áp dụng.
\loigiai{
Cần nêu được: Độ phóng xạ $H=\lambda N$ và cũng giảm theo quy luật mũ với cùng chu kì bán rã. Khi vận dụng phải xác định đúng dữ kiện, điều kiện áp dụng, hệ đơn vị và kiểm tra tính hợp lí của kết quả. Nhận định “Độ phóng xạ không phụ thuộc số hạt chưa phân rã.” sai vì trái với kiến thức trên.
}
\end{bt}

% ===== Câu 99 =====
% ID: L12C4B23-25-TLN
% Mức: TH
\begin{ex}
% ID: L12C4B23-25-TLN
% Mức: TH
Trong bốn phát biểu sau về tính độ phóng xạ và khai thác đồ thị phóng xạ, có bao nhiêu phát biểu đúng? Chỉ ghi một số nguyên. (1) Độ phóng xạ $H=\lambda N$ và cũng giảm theo quy luật mũ với cùng chu kì bán rã. (2) Độ phóng xạ không phụ thuộc số hạt chưa phân rã. (3) Cần kiểm tra điều kiện áp dụng trước khi tính toán. (4) Có thể bỏ qua đơn vị của kết quả.
\shortans{2}
\loigiai{
Đối chiếu từng phát biểu với kiến thức cốt lõi: Độ phóng xạ $H=\lambda N$ và cũng giảm theo quy luật mũ với cùng chu kì bán rã. Có 2 phát biểu đúng.
}
\end{ex}

% ===== Câu 100 =====
% ID: L12C4B23-25-TN
% Mức: TH
\dangbt{Vận dụng định luật phóng xạ theo số hạt và khối lượng}
\begin{ex}
% ID: L12C4B23-25-TN
% Mức: TH
Khi phân tích, phát biểu nào đúng về vận dụng định luật phóng xạ theo số hạt và khối lượng?
\choice
{\True Số hạt chưa phân rã và khối lượng chất phóng xạ giảm theo $N=N_0 2^{-t/T}$ và $m=m_0 2^{-t/T}$.}
{Sau mỗi chu kì bán rã, lượng chất chưa phân rã giảm một lượng cố định như nhau.}
{Có thể bỏ qua điều kiện áp dụng và đơn vị trong mọi trường hợp.}
{Kết quả của bài toán không phụ thuộc dữ kiện đã cho.}
\loigiai{
Số hạt chưa phân rã và khối lượng chất phóng xạ giảm theo $N=N_0 2^{-t/T}$ và $m=m_0 2^{-t/T}$. Vì vậy phương án $A$ đúng; các phương án còn lại trái với kiến thức hoặc điều kiện áp dụng.
}
\end{ex}

% ===== Câu 101 =====
% ID: L12C4B23-26-DS
% Mức: TH
\dangbt{Ứng dụng, ảnh hưởng và an toàn phóng xạ}
\begin{ex}
% ID: L12C4B23-26-DS
% Mức: TH
Xét ứng dụng, ảnh hưởng và an toàn phóng xạ (tình huống 1). Hãy xác định tính đúng, sai của bốn phát biểu sau.
\choiceTF
{\True Phóng xạ được dùng trong y học, công nghiệp và định tuổi; an toàn dựa trên giảm thời gian, tăng khoảng cách và che chắn phù hợp.}
{Có thể tiếp xúc nguồn phóng xạ lâu nếu nguồn có kích thước nhỏ.}
{\True Khi giải cần đọc đúng dữ kiện, dùng hệ đơn vị thống nhất và kiểm tra điều kiện áp dụng.}
{Có thể bỏ qua dữ kiện, đơn vị và ý nghĩa vật lí của kết quả.}
\loigiai{
\begin{itemchoice}
\item (Đúng) Phóng xạ được dùng trong y học, công nghiệp và định tuổi; an toàn dựa trên giảm thời gian, tăng khoảng cách và che chắn phù hợp. (Vì): Phóng xạ được dùng trong y học, công nghiệp và định tuổi; an toàn dựa trên giảm thời gian, tăng khoảng cách và che chắn phù hợp. b) Sai: Có thể tiếp xúc nguồn phóng xạ lâu nếu nguồn có kích thước nhỏ. c) Đúng: Khi giải cần đọc đúng dữ kiện, dùng hệ đơn vị thống nhất và kiểm tra điều kiện áp dụng. d) Sai: Có thể bỏ qua dữ kiện, đơn vị và ý nghĩa vật lí của kết quả. Kiến thức cốt lõi: Phóng xạ được dùng trong y học, công nghiệp và định tuổi; an toàn dựa trên giảm thời gian, tăng khoảng cách và che chắn phù hợp
\item (Sai) Có thể tiếp xúc nguồn phóng xạ lâu nếu nguồn có kích thước nhỏ.
\item (Đúng) Khi giải cần đọc đúng dữ kiện, dùng hệ đơn vị thống nhất và kiểm tra điều kiện áp dụng.
\item (Sai) Có thể bỏ qua dữ kiện, đơn vị và ý nghĩa vật lí của kết quả.
\end{itemchoice}
}
\end{ex}

% ===== Câu 102 =====
% ID: L12C4B23-26-TL
% Mức: TH
\dangbt{Tính độ phóng xạ và khai thác đồ thị phóng xạ}
\begin{bt}
% ID: L12C4B23-26-TL
% Mức: TH
Giải thích vì sao nhận định sau đúng: “Độ phóng xạ $H=\lambda N$ và cũng giảm theo quy luật mũ với cùng chu kì bán rã.”
\loigiai{
Cần nêu được: Độ phóng xạ $H=\lambda N$ và cũng giảm theo quy luật mũ với cùng chu kì bán rã. Khi vận dụng phải xác định đúng dữ kiện, điều kiện áp dụng, hệ đơn vị và kiểm tra tính hợp lí của kết quả. Nhận định “Độ phóng xạ không phụ thuộc số hạt chưa phân rã.” sai vì trái với kiến thức trên.
}
\end{bt}

% ===== Câu 103 =====
% ID: L12C4B23-26-TLN
% Mức: TH
\begin{ex}
% ID: L12C4B23-26-TLN
% Mức: TH
Trong bốn phát biểu sau về tính độ phóng xạ và khai thác đồ thị phóng xạ, có bao nhiêu phát biểu đúng? Chỉ ghi một số nguyên. (1) Độ phóng xạ không phụ thuộc số hạt chưa phân rã. (2) Độ phóng xạ $H=\lambda N$ và cũng giảm theo quy luật mũ với cùng chu kì bán rã. (3) Kết quả phải phù hợp ý nghĩa vật lí. (4) Mọi đại lượng đều có cùng bản chất.
\shortans{2}
\loigiai{
Đối chiếu từng phát biểu với kiến thức cốt lõi: Độ phóng xạ $H=\lambda N$ và cũng giảm theo quy luật mũ với cùng chu kì bán rã. Có 2 phát biểu đúng.
}
\end{ex}

% ===== Câu 104 =====
% ID: L12C4B23-26-TN
% Mức: TH
\dangbt{Vận dụng định luật phóng xạ theo số hạt và khối lượng}
\begin{ex}
% ID: L12C4B23-26-TN
% Mức: TH
Kết luận nào của học sinh là đúng về vận dụng định luật phóng xạ theo số hạt và khối lượng?
\choice
{Sau mỗi chu kì bán rã, lượng chất chưa phân rã giảm một lượng cố định như nhau.}
{\True Số hạt chưa phân rã và khối lượng chất phóng xạ giảm theo $N=N_0 2^{-t/T}$ và $m=m_0 2^{-t/T}$.}
{Có thể bỏ qua điều kiện áp dụng và đơn vị trong mọi trường hợp.}
{Kết quả của bài toán không phụ thuộc dữ kiện đã cho.}
\loigiai{
Số hạt chưa phân rã và khối lượng chất phóng xạ giảm theo $N=N_0 2^{-t/T}$ và $m=m_0 2^{-t/T}$. Vì vậy phương án $B$ đúng; các phương án còn lại trái với kiến thức hoặc điều kiện áp dụng.
}
\end{ex}

% ===== Câu 105 =====
% ID: L12C4B23-27-DS
% Mức: TH
\dangbt{Ứng dụng, ảnh hưởng và an toàn phóng xạ}
\begin{ex}
% ID: L12C4B23-27-DS
% Mức: TH
Xét ứng dụng, ảnh hưởng và an toàn phóng xạ (tình huống 2). Hãy xác định tính đúng, sai của bốn phát biểu sau.
\choiceTF
{Có thể tiếp xúc nguồn phóng xạ lâu nếu nguồn có kích thước nhỏ.}
{\True Phóng xạ được dùng trong y học, công nghiệp và định tuổi; an toàn dựa trên giảm thời gian, tăng khoảng cách và che chắn phù hợp.}
{Có thể bỏ qua dữ kiện, đơn vị và ý nghĩa vật lí của kết quả.}
{\True Khi giải cần đọc đúng dữ kiện, dùng hệ đơn vị thống nhất và kiểm tra điều kiện áp dụng.}
\loigiai{
\begin{itemchoice}
\item (Sai) Có thể tiếp xúc nguồn phóng xạ lâu nếu nguồn có kích thước nhỏ. (Vì): Có thể tiếp xúc nguồn phóng xạ lâu nếu nguồn có kích thước nhỏ. b) Đúng: Phóng xạ được dùng trong y học, công nghiệp và định tuổi; an toàn dựa trên giảm thời gian, tăng khoảng cách và che chắn phù hợp. c) Sai: Có thể bỏ qua dữ kiện, đơn vị và ý nghĩa vật lí của kết quả. d) Đúng: Khi giải cần đọc đúng dữ kiện, dùng hệ đơn vị thống nhất và kiểm tra điều kiện áp dụng. Kiến thức cốt lõi: Phóng xạ được dùng trong y học, công nghiệp và định tuổi; an toàn dựa trên giảm thời gian, tăng khoảng cách và che chắn phù hợp
\item (Đúng) Phóng xạ được dùng trong y học, công nghiệp và định tuổi; an toàn dựa trên giảm thời gian, tăng khoảng cách và che chắn phù hợp.
\item (Sai) Có thể bỏ qua dữ kiện, đơn vị và ý nghĩa vật lí của kết quả.
\item (Đúng) Khi giải cần đọc đúng dữ kiện, dùng hệ đơn vị thống nhất và kiểm tra điều kiện áp dụng.
\end{itemchoice}
}
\end{ex}

% ===== Câu 106 =====
% ID: L12C4B23-27-TL
% Mức: VD
\dangbt{Tính độ phóng xạ và khai thác đồ thị phóng xạ}
\begin{bt}
% ID: L12C4B23-27-TL
% Mức: VD
Phân tích sai lầm trong nhận định: “Độ phóng xạ không phụ thuộc số hạt chưa phân rã.”
\loigiai{
Cần nêu được: Độ phóng xạ $H=\lambda N$ và cũng giảm theo quy luật mũ với cùng chu kì bán rã. Khi vận dụng phải xác định đúng dữ kiện, điều kiện áp dụng, hệ đơn vị và kiểm tra tính hợp lí của kết quả. Nhận định “Độ phóng xạ không phụ thuộc số hạt chưa phân rã.” sai vì trái với kiến thức trên.
}
\end{bt}

% ===== Câu 107 =====
% ID: L12C4B23-27-TLN
% Mức: TH
\begin{ex}
% ID: L12C4B23-27-TLN
% Mức: TH
Trong bốn phát biểu sau về tính độ phóng xạ và khai thác đồ thị phóng xạ, có bao nhiêu phát biểu đúng? Chỉ ghi một số nguyên. (1) Cần đổi các đại lượng về hệ đơn vị thống nhất. (2) Độ phóng xạ $H=\lambda N$ và cũng giảm theo quy luật mũ với cùng chu kì bán rã. (3) Độ phóng xạ không phụ thuộc số hạt chưa phân rã. (4) Không cần kiểm tra dữ kiện.
\shortans{2}
\loigiai{
Đối chiếu từng phát biểu với kiến thức cốt lõi: Độ phóng xạ $H=\lambda N$ và cũng giảm theo quy luật mũ với cùng chu kì bán rã. Có 2 phát biểu đúng.
}
\end{ex}

% ===== Câu 108 =====
% ID: L12C4B23-27-TN
% Mức: TH
\dangbt{Vận dụng định luật phóng xạ theo số hạt và khối lượng}
\begin{ex}
% ID: L12C4B23-27-TN
% Mức: TH
Chọn phát biểu không trái với kiến thức về vận dụng định luật phóng xạ theo số hạt và khối lượng?
\choice
{Sau mỗi chu kì bán rã, lượng chất chưa phân rã giảm một lượng cố định như nhau.}
{Có thể bỏ qua điều kiện áp dụng và đơn vị trong mọi trường hợp.}
{\True Số hạt chưa phân rã và khối lượng chất phóng xạ giảm theo $N=N_0 2^{-t/T}$ và $m=m_0 2^{-t/T}$.}
{Kết quả của bài toán không phụ thuộc dữ kiện đã cho.}
\loigiai{
Số hạt chưa phân rã và khối lượng chất phóng xạ giảm theo $N=N_0 2^{-t/T}$ và $m=m_0 2^{-t/T}$. Vì vậy phương án $C$ đúng; các phương án còn lại trái với kiến thức hoặc điều kiện áp dụng.
}
\end{ex}

% ===== Câu 109 =====
% ID: L12C4B23-28-DS
% Mức: VD
\dangbt{Ứng dụng, ảnh hưởng và an toàn phóng xạ}
\begin{ex}
% ID: L12C4B23-28-DS
% Mức: VD
Xét ứng dụng, ảnh hưởng và an toàn phóng xạ (tình huống 3). Hãy xác định tính đúng, sai của bốn phát biểu sau.
\choiceTF
{\True Phóng xạ được dùng trong y học, công nghiệp và định tuổi; an toàn dựa trên giảm thời gian, tăng khoảng cách và che chắn phù hợp.}
{Có thể tiếp xúc nguồn phóng xạ lâu nếu nguồn có kích thước nhỏ.}
{\True Khi giải cần đọc đúng dữ kiện, dùng hệ đơn vị thống nhất và kiểm tra điều kiện áp dụng.}
{Có thể bỏ qua dữ kiện, đơn vị và ý nghĩa vật lí của kết quả.}
\loigiai{
\begin{itemchoice}
\item (Đúng) Phóng xạ được dùng trong y học, công nghiệp và định tuổi; an toàn dựa trên giảm thời gian, tăng khoảng cách và che chắn phù hợp. (Vì): Phóng xạ được dùng trong y học, công nghiệp và định tuổi; an toàn dựa trên giảm thời gian, tăng khoảng cách và che chắn phù hợp. b) Sai: Có thể tiếp xúc nguồn phóng xạ lâu nếu nguồn có kích thước nhỏ. c) Đúng: Khi giải cần đọc đúng dữ kiện, dùng hệ đơn vị thống nhất và kiểm tra điều kiện áp dụng. d) Sai: Có thể bỏ qua dữ kiện, đơn vị và ý nghĩa vật lí của kết quả. Kiến thức cốt lõi: Phóng xạ được dùng trong y học, công nghiệp và định tuổi; an toàn dựa trên giảm thời gian, tăng khoảng cách và che chắn phù hợp
\item (Sai) Có thể tiếp xúc nguồn phóng xạ lâu nếu nguồn có kích thước nhỏ.
\item (Đúng) Khi giải cần đọc đúng dữ kiện, dùng hệ đơn vị thống nhất và kiểm tra điều kiện áp dụng.
\item (Sai) Có thể bỏ qua dữ kiện, đơn vị và ý nghĩa vật lí của kết quả.
\end{itemchoice}
}
\end{ex}

% ===== Câu 110 =====
% ID: L12C4B23-28-TL
% Mức: VD
\dangbt{Tính độ phóng xạ và khai thác đồ thị phóng xạ}
\begin{bt}
% ID: L12C4B23-28-TL
% Mức: VD
Đề xuất các bước giải một bài toán thuộc dạng tính độ phóng xạ và khai thác đồ thị phóng xạ.
\loigiai{
Cần nêu được: Độ phóng xạ $H=\lambda N$ và cũng giảm theo quy luật mũ với cùng chu kì bán rã. Khi vận dụng phải xác định đúng dữ kiện, điều kiện áp dụng, hệ đơn vị và kiểm tra tính hợp lí của kết quả. Nhận định “Độ phóng xạ không phụ thuộc số hạt chưa phân rã.” sai vì trái với kiến thức trên.
}
\end{bt}

% ===== Câu 111 =====
% ID: L12C4B23-28-TLN
% Mức: VD
\begin{ex}
% ID: L12C4B23-28-TLN
% Mức: VD
Trong bốn phát biểu sau về tính độ phóng xạ và khai thác đồ thị phóng xạ, có bao nhiêu phát biểu đúng? Chỉ ghi một số nguyên. (1) Độ phóng xạ $H=\lambda N$ và cũng giảm theo quy luật mũ với cùng chu kì bán rã. (2) Cần đối chiếu kết quả với điều kiện bài toán. (3) Độ phóng xạ không phụ thuộc số hạt chưa phân rã. (4) Mọi trường hợp đều cho cùng một kết quả.
\shortans{2}
\loigiai{
Đối chiếu từng phát biểu với kiến thức cốt lõi: Độ phóng xạ $H=\lambda N$ và cũng giảm theo quy luật mũ với cùng chu kì bán rã. Có 2 phát biểu đúng.
}
\end{ex}

% ===== Câu 112 =====
% ID: L12C4B23-28-TN
% Mức: VD
\dangbt{Vận dụng định luật phóng xạ theo số hạt và khối lượng}
\begin{ex}
% ID: L12C4B23-28-TN
% Mức: VD
Cơ sở Vật lí đúng để giải bài là gì về vận dụng định luật phóng xạ theo số hạt và khối lượng?
\choice
{Sau mỗi chu kì bán rã, lượng chất chưa phân rã giảm một lượng cố định như nhau.}
{Có thể bỏ qua điều kiện áp dụng và đơn vị trong mọi trường hợp.}
{Kết quả của bài toán không phụ thuộc dữ kiện đã cho.}
{\True Số hạt chưa phân rã và khối lượng chất phóng xạ giảm theo $N=N_0 2^{-t/T}$ và $m=m_0 2^{-t/T}$.}
\loigiai{
Số hạt chưa phân rã và khối lượng chất phóng xạ giảm theo $N=N_0 2^{-t/T}$ và $m=m_0 2^{-t/T}$. Vì vậy phương án $D$ đúng; các phương án còn lại trái với kiến thức hoặc điều kiện áp dụng.
}
\end{ex}

% ===== Câu 113 =====
% ID: L12C4B23-29-DS
% Mức: VD
\dangbt{Ứng dụng, ảnh hưởng và an toàn phóng xạ}
\begin{ex}
% ID: L12C4B23-29-DS
% Mức: VD
Xét ứng dụng, ảnh hưởng và an toàn phóng xạ (tình huống 4). Hãy xác định tính đúng, sai của bốn phát biểu sau.
\choiceTF
{Có thể tiếp xúc nguồn phóng xạ lâu nếu nguồn có kích thước nhỏ.}
{\True Phóng xạ được dùng trong y học, công nghiệp và định tuổi; an toàn dựa trên giảm thời gian, tăng khoảng cách và che chắn phù hợp.}
{Có thể bỏ qua dữ kiện, đơn vị và ý nghĩa vật lí của kết quả.}
{\True Khi giải cần đọc đúng dữ kiện, dùng hệ đơn vị thống nhất và kiểm tra điều kiện áp dụng.}
\loigiai{
\begin{itemchoice}
\item (Sai) Có thể tiếp xúc nguồn phóng xạ lâu nếu nguồn có kích thước nhỏ. (Vì): Có thể tiếp xúc nguồn phóng xạ lâu nếu nguồn có kích thước nhỏ. b) Đúng: Phóng xạ được dùng trong y học, công nghiệp và định tuổi; an toàn dựa trên giảm thời gian, tăng khoảng cách và che chắn phù hợp. c) Sai: Có thể bỏ qua dữ kiện, đơn vị và ý nghĩa vật lí của kết quả. d) Đúng: Khi giải cần đọc đúng dữ kiện, dùng hệ đơn vị thống nhất và kiểm tra điều kiện áp dụng. Kiến thức cốt lõi: Phóng xạ được dùng trong y học, công nghiệp và định tuổi; an toàn dựa trên giảm thời gian, tăng khoảng cách và che chắn phù hợp
\item (Đúng) Phóng xạ được dùng trong y học, công nghiệp và định tuổi; an toàn dựa trên giảm thời gian, tăng khoảng cách và che chắn phù hợp.
\item (Sai) Có thể bỏ qua dữ kiện, đơn vị và ý nghĩa vật lí của kết quả.
\item (Đúng) Khi giải cần đọc đúng dữ kiện, dùng hệ đơn vị thống nhất và kiểm tra điều kiện áp dụng.
\end{itemchoice}
}
\end{ex}

% ===== Câu 114 =====
% ID: L12C4B23-29-TL
% Mức: VD
\dangbt{Tính độ phóng xạ và khai thác đồ thị phóng xạ}
\begin{bt}
% ID: L12C4B23-29-TL
% Mức: VD
Nêu một tình huống thực tiễn liên quan đến tính độ phóng xạ và khai thác đồ thị phóng xạ và giải thích bằng kiến thức Vật lí.
\loigiai{
Cần nêu được: Độ phóng xạ $H=\lambda N$ và cũng giảm theo quy luật mũ với cùng chu kì bán rã. Khi vận dụng phải xác định đúng dữ kiện, điều kiện áp dụng, hệ đơn vị và kiểm tra tính hợp lí của kết quả. Nhận định “Độ phóng xạ không phụ thuộc số hạt chưa phân rã.” sai vì trái với kiến thức trên.
}
\end{bt}

% ===== Câu 115 =====
% ID: L12C4B23-29-TLN
% Mức: VD
\begin{ex}
% ID: L12C4B23-29-TLN
% Mức: VD
Trong bốn phát biểu sau về tính độ phóng xạ và khai thác đồ thị phóng xạ, có bao nhiêu phát biểu đúng? Chỉ ghi một số nguyên. (1) Độ phóng xạ không phụ thuộc số hạt chưa phân rã. (2) Cần đọc đúng dữ kiện và giả thiết. (3) Độ phóng xạ $H=\lambda N$ và cũng giảm theo quy luật mũ với cùng chu kì bán rã. (4) Có thể thay số trước khi chọn công thức.
\shortans{2}
\loigiai{
Đối chiếu từng phát biểu với kiến thức cốt lõi: Độ phóng xạ $H=\lambda N$ và cũng giảm theo quy luật mũ với cùng chu kì bán rã. Có 2 phát biểu đúng.
}
\end{ex}

% ===== Câu 116 =====
% ID: L12C4B23-29-TN
% Mức: VD
\dangbt{Vận dụng định luật phóng xạ theo số hạt và khối lượng}
\begin{ex}
% ID: L12C4B23-29-TN
% Mức: VD
Trong một bài thực hành, nhận xét nào đúng về vận dụng định luật phóng xạ theo số hạt và khối lượng?
\choice
{\True Số hạt chưa phân rã và khối lượng chất phóng xạ giảm theo $N=N_0 2^{-t/T}$ và $m=m_0 2^{-t/T}$.}
{Sau mỗi chu kì bán rã, lượng chất chưa phân rã giảm một lượng cố định như nhau.}
{Có thể bỏ qua điều kiện áp dụng và đơn vị trong mọi trường hợp.}
{Kết quả của bài toán không phụ thuộc dữ kiện đã cho.}
\loigiai{
Số hạt chưa phân rã và khối lượng chất phóng xạ giảm theo $N=N_0 2^{-t/T}$ và $m=m_0 2^{-t/T}$. Vì vậy phương án $A$ đúng; các phương án còn lại trái với kiến thức hoặc điều kiện áp dụng.
}
\end{ex}

% ===== Câu 117 =====
% ID: L12C4B23-30-DS
% Mức: VD
\dangbt{Ứng dụng, ảnh hưởng và an toàn phóng xạ}
\begin{ex}
% ID: L12C4B23-30-DS
% Mức: VD
Xét ứng dụng, ảnh hưởng và an toàn phóng xạ (tình huống 5). Hãy xác định tính đúng, sai của bốn phát biểu sau.
\choiceTF
{\True Phóng xạ được dùng trong y học, công nghiệp và định tuổi; an toàn dựa trên giảm thời gian, tăng khoảng cách và che chắn phù hợp.}
{Có thể bỏ qua dữ kiện, đơn vị và ý nghĩa vật lí của kết quả.}
{Có thể tiếp xúc nguồn phóng xạ lâu nếu nguồn có kích thước nhỏ.}
{Có thể bỏ qua dữ kiện, đơn vị và ý nghĩa vật lí của kết quả.}
\loigiai{
\begin{itemchoice}
\item (Đúng) Phóng xạ được dùng trong y học, công nghiệp và định tuổi; an toàn dựa trên giảm thời gian, tăng khoảng cách và che chắn phù hợp. (Vì): Phóng xạ được dùng trong y học, công nghiệp và định tuổi; an toàn dựa trên giảm thời gian, tăng khoảng cách và che chắn phù hợp. b) Sai: Có thể bỏ qua dữ kiện, đơn vị và ý nghĩa vật lí của kết quả. c) Sai: Có thể tiếp xúc nguồn phóng xạ lâu nếu nguồn có kích thước nhỏ. d) Sai: Có thể bỏ qua dữ kiện, đơn vị và ý nghĩa vật lí của kết quả. Kiến thức cốt lõi: Phóng xạ được dùng trong y học, công nghiệp và định tuổi; an toàn dựa trên giảm thời gian, tăng khoảng cách và che chắn phù hợp
\item (Sai) Có thể bỏ qua dữ kiện, đơn vị và ý nghĩa vật lí của kết quả.
\item (Sai) Có thể tiếp xúc nguồn phóng xạ lâu nếu nguồn có kích thước nhỏ.
\item (Sai) Có thể bỏ qua dữ kiện, đơn vị và ý nghĩa vật lí của kết quả.
\end{itemchoice}
}
\end{ex}

% ===== Câu 118 =====
% ID: L12C4B23-30-TL
% Mức: VDC
\dangbt{Tính độ phóng xạ và khai thác đồ thị phóng xạ}
\begin{bt}
% ID: L12C4B23-30-TL
% Mức: VDC
So sánh nhận định đúng “Độ phóng xạ $H=\lambda N$ và cũng giảm theo quy luật mũ với cùng chu kì bán rã.” với nhận định sai “Độ phóng xạ không phụ thuộc số hạt chưa phân rã.”; chỉ rõ căn cứ Vật lí.
\loigiai{
Cần nêu được: Độ phóng xạ $H=\lambda N$ và cũng giảm theo quy luật mũ với cùng chu kì bán rã. Khi vận dụng phải xác định đúng dữ kiện, điều kiện áp dụng, hệ đơn vị và kiểm tra tính hợp lí của kết quả. Nhận định “Độ phóng xạ không phụ thuộc số hạt chưa phân rã.” sai vì trái với kiến thức trên.
}
\end{bt}

% ===== Câu 119 =====
% ID: L12C4B23-30-TLN
% Mức: VDC
\begin{ex}
% ID: L12C4B23-30-TLN
% Mức: VDC
Trong bốn phát biểu sau về tính độ phóng xạ và khai thác đồ thị phóng xạ, có bao nhiêu phát biểu đúng? Chỉ ghi một số nguyên. (1) Kết quả cần có ý nghĩa vật lí. (2) Độ phóng xạ không phụ thuộc số hạt chưa phân rã. (3) Phải dùng đúng định luật cho tình huống. (4) Độ phóng xạ $H=\lambda N$ và cũng giảm theo quy luật mũ với cùng chu kì bán rã.
\shortans{3}
\loigiai{
Đối chiếu từng phát biểu với kiến thức cốt lõi: Độ phóng xạ $H=\lambda N$ và cũng giảm theo quy luật mũ với cùng chu kì bán rã. Có 3 phát biểu đúng.
}
\end{ex}

% ===== Câu 120 =====
% ID: L12C4B23-30-TN
% Mức: VD
\dangbt{Vận dụng định luật phóng xạ theo số hạt và khối lượng}
\begin{ex}
% ID: L12C4B23-30-TN
% Mức: VD
Kết luận đúng khi vận dụng là về vận dụng định luật phóng xạ theo số hạt và khối lượng?
\choice
{Sau mỗi chu kì bán rã, lượng chất chưa phân rã giảm một lượng cố định như nhau.}
{\True Số hạt chưa phân rã và khối lượng chất phóng xạ giảm theo $N=N_0 2^{-t/T}$ và $m=m_0 2^{-t/T}$.}
{Có thể bỏ qua điều kiện áp dụng và đơn vị trong mọi trường hợp.}
{Kết quả của bài toán không phụ thuộc dữ kiện đã cho.}
\loigiai{
Số hạt chưa phân rã và khối lượng chất phóng xạ giảm theo $N=N_0 2^{-t/T}$ và $m=m_0 2^{-t/T}$. Vì vậy phương án $B$ đúng; các phương án còn lại trái với kiến thức hoặc điều kiện áp dụng.
}
\end{ex}

% ===== Câu 121 =====
% ID: L12C4B23-31-TL
% Mức: TH
\dangbt{Ứng dụng, ảnh hưởng và an toàn phóng xạ}
\begin{bt}
% ID: L12C4B23-31-TL
% Mức: TH
Trình bày kiến thức cốt lõi của ứng dụng, ảnh hưởng và an toàn phóng xạ và nêu điều kiện áp dụng.
\loigiai{
Cần nêu được: Phóng xạ được dùng trong y học, công nghiệp và định tuổi; an toàn dựa trên giảm thời gian, tăng khoảng cách và che chắn phù hợp. Khi vận dụng phải xác định đúng dữ kiện, điều kiện áp dụng, hệ đơn vị và kiểm tra tính hợp lí của kết quả. Nhận định “Có thể tiếp xúc nguồn phóng xạ lâu nếu nguồn có kích thước nhỏ.” sai vì trái với kiến thức trên.
}
\end{bt}

% ===== Câu 122 =====
% ID: L12C4B23-31-TLN
% Mức: TH
\begin{ex}
% ID: L12C4B23-31-TLN
% Mức: TH
Trong bốn phát biểu sau về ứng dụng, ảnh hưởng và an toàn phóng xạ, có bao nhiêu phát biểu đúng? Chỉ ghi một số nguyên. (1) Phóng xạ được dùng trong y học, công nghiệp và định tuổi; an toàn dựa trên giảm thời gian, tăng khoảng cách và che chắn phù hợp. (2) Có thể tiếp xúc nguồn phóng xạ lâu nếu nguồn có kích thước nhỏ. (3) Cần kiểm tra điều kiện áp dụng trước khi tính toán. (4) Có thể bỏ qua đơn vị của kết quả.
\shortans{2}
\loigiai{
Đối chiếu từng phát biểu với kiến thức cốt lõi: Phóng xạ được dùng trong y học, công nghiệp và định tuổi; an toàn dựa trên giảm thời gian, tăng khoảng cách và che chắn phù hợp. Có 2 phát biểu đúng.
}
\end{ex}

% ===== Câu 123 =====
% ID: L12C4B23-31-TN
% Mức: NB
\dangbt{Tính chu kì bán rã, hằng số phóng xạ và tuổi mẫu vật}
\begin{ex}
% ID: L12C4B23-31-TN
% Mức: NB
Phát biểu nào sau đây đúng về tính chu kì bán rã, hằng số phóng xạ và tuổi mẫu vật?
\choice
{\True Hằng số phóng xạ $\lambda=\ln2/T$; từ tỉ số lượng còn lại có thể suy ra thời gian hay tuổi mẫu.}
{Chu kì bán rã càng lớn thì hằng số phóng xạ càng lớn.}
{Có thể bỏ qua điều kiện áp dụng và đơn vị trong mọi trường hợp.}
{Kết quả của bài toán không phụ thuộc dữ kiện đã cho.}
\loigiai{
Hằng số phóng xạ $\lambda=\ln2/T$; từ tỉ số lượng còn lại có thể suy ra thời gian hay tuổi mẫu. Vì vậy phương án $A$ đúng; các phương án còn lại trái với kiến thức hoặc điều kiện áp dụng.
}
\end{ex}

% ===== Câu 124 =====
% ID: L12C4B23-32-TL
% Mức: TH
\dangbt{Ứng dụng, ảnh hưởng và an toàn phóng xạ}
\begin{bt}
% ID: L12C4B23-32-TL
% Mức: TH
Giải thích vì sao nhận định sau đúng: “Phóng xạ được dùng trong y học, công nghiệp và định tuổi; an toàn dựa trên giảm thời gian, tăng khoảng cách và che chắn phù hợp.”
\loigiai{
Cần nêu được: Phóng xạ được dùng trong y học, công nghiệp và định tuổi; an toàn dựa trên giảm thời gian, tăng khoảng cách và che chắn phù hợp. Khi vận dụng phải xác định đúng dữ kiện, điều kiện áp dụng, hệ đơn vị và kiểm tra tính hợp lí của kết quả. Nhận định “Có thể tiếp xúc nguồn phóng xạ lâu nếu nguồn có kích thước nhỏ.” sai vì trái với kiến thức trên.
}
\end{bt}

% ===== Câu 125 =====
% ID: L12C4B23-32-TLN
% Mức: TH
\begin{ex}
% ID: L12C4B23-32-TLN
% Mức: TH
Trong bốn phát biểu sau về ứng dụng, ảnh hưởng và an toàn phóng xạ, có bao nhiêu phát biểu đúng? Chỉ ghi một số nguyên. (1) Có thể tiếp xúc nguồn phóng xạ lâu nếu nguồn có kích thước nhỏ. (2) Phóng xạ được dùng trong y học, công nghiệp và định tuổi; an toàn dựa trên giảm thời gian, tăng khoảng cách và che chắn phù hợp. (3) Kết quả phải phù hợp ý nghĩa vật lí. (4) Mọi đại lượng đều có cùng bản chất.
\shortans{2}
\loigiai{
Đối chiếu từng phát biểu với kiến thức cốt lõi: Phóng xạ được dùng trong y học, công nghiệp và định tuổi; an toàn dựa trên giảm thời gian, tăng khoảng cách và che chắn phù hợp. Có 2 phát biểu đúng.
}
\end{ex}

% ===== Câu 126 =====
% ID: L12C4B23-32-TN
% Mức: NB
\dangbt{Tính chu kì bán rã, hằng số phóng xạ và tuổi mẫu vật}
\begin{ex}
% ID: L12C4B23-32-TN
% Mức: NB
Nhận định nào phù hợp với kiến thức Vật lí về tính chu kì bán rã, hằng số phóng xạ và tuổi mẫu vật?
\choice
{Chu kì bán rã càng lớn thì hằng số phóng xạ càng lớn.}
{\True Hằng số phóng xạ $\lambda=\ln2/T$; từ tỉ số lượng còn lại có thể suy ra thời gian hay tuổi mẫu.}
{Có thể bỏ qua điều kiện áp dụng và đơn vị trong mọi trường hợp.}
{Kết quả của bài toán không phụ thuộc dữ kiện đã cho.}
\loigiai{
Hằng số phóng xạ $\lambda=\ln2/T$; từ tỉ số lượng còn lại có thể suy ra thời gian hay tuổi mẫu. Vì vậy phương án $B$ đúng; các phương án còn lại trái với kiến thức hoặc điều kiện áp dụng.
}
\end{ex}

% ===== Câu 127 =====
% ID: L12C4B23-33-TL
% Mức: VD
\dangbt{Ứng dụng, ảnh hưởng và an toàn phóng xạ}
\begin{bt}
% ID: L12C4B23-33-TL
% Mức: VD
Phân tích sai lầm trong nhận định: “Có thể tiếp xúc nguồn phóng xạ lâu nếu nguồn có kích thước nhỏ.”
\loigiai{
Cần nêu được: Phóng xạ được dùng trong y học, công nghiệp và định tuổi; an toàn dựa trên giảm thời gian, tăng khoảng cách và che chắn phù hợp. Khi vận dụng phải xác định đúng dữ kiện, điều kiện áp dụng, hệ đơn vị và kiểm tra tính hợp lí của kết quả. Nhận định “Có thể tiếp xúc nguồn phóng xạ lâu nếu nguồn có kích thước nhỏ.” sai vì trái với kiến thức trên.
}
\end{bt}

% ===== Câu 128 =====
% ID: L12C4B23-33-TLN
% Mức: TH
\begin{ex}
% ID: L12C4B23-33-TLN
% Mức: TH
Trong bốn phát biểu sau về ứng dụng, ảnh hưởng và an toàn phóng xạ, có bao nhiêu phát biểu đúng? Chỉ ghi một số nguyên. (1) Cần đổi các đại lượng về hệ đơn vị thống nhất. (2) Phóng xạ được dùng trong y học, công nghiệp và định tuổi; an toàn dựa trên giảm thời gian, tăng khoảng cách và che chắn phù hợp. (3) Có thể tiếp xúc nguồn phóng xạ lâu nếu nguồn có kích thước nhỏ. (4) Không cần kiểm tra dữ kiện.
\shortans{2}
\loigiai{
Đối chiếu từng phát biểu với kiến thức cốt lõi: Phóng xạ được dùng trong y học, công nghiệp và định tuổi; an toàn dựa trên giảm thời gian, tăng khoảng cách và che chắn phù hợp. Có 2 phát biểu đúng.
}
\end{ex}

% ===== Câu 129 =====
% ID: L12C4B23-33-TN
% Mức: NB
\dangbt{Tính chu kì bán rã, hằng số phóng xạ và tuổi mẫu vật}
\begin{ex}
% ID: L12C4B23-33-TN
% Mức: NB
Chọn kết luận chính xác về tính chu kì bán rã, hằng số phóng xạ và tuổi mẫu vật?
\choice
{Chu kì bán rã càng lớn thì hằng số phóng xạ càng lớn.}
{Có thể bỏ qua điều kiện áp dụng và đơn vị trong mọi trường hợp.}
{\True Hằng số phóng xạ $\lambda=\ln2/T$; từ tỉ số lượng còn lại có thể suy ra thời gian hay tuổi mẫu.}
{Kết quả của bài toán không phụ thuộc dữ kiện đã cho.}
\loigiai{
Hằng số phóng xạ $\lambda=\ln2/T$; từ tỉ số lượng còn lại có thể suy ra thời gian hay tuổi mẫu. Vì vậy phương án $C$ đúng; các phương án còn lại trái với kiến thức hoặc điều kiện áp dụng.
}
\end{ex}

% ===== Câu 130 =====
% ID: L12C4B23-34-TL
% Mức: VD
\dangbt{Ứng dụng, ảnh hưởng và an toàn phóng xạ}
\begin{bt}
% ID: L12C4B23-34-TL
% Mức: VD
Đề xuất các bước giải một bài toán thuộc dạng ứng dụng, ảnh hưởng và an toàn phóng xạ.
\loigiai{
Cần nêu được: Phóng xạ được dùng trong y học, công nghiệp và định tuổi; an toàn dựa trên giảm thời gian, tăng khoảng cách và che chắn phù hợp. Khi vận dụng phải xác định đúng dữ kiện, điều kiện áp dụng, hệ đơn vị và kiểm tra tính hợp lí của kết quả. Nhận định “Có thể tiếp xúc nguồn phóng xạ lâu nếu nguồn có kích thước nhỏ.” sai vì trái với kiến thức trên.
}
\end{bt}

% ===== Câu 131 =====
% ID: L12C4B23-34-TLN
% Mức: VD
\begin{ex}
% ID: L12C4B23-34-TLN
% Mức: VD
Trong bốn phát biểu sau về ứng dụng, ảnh hưởng và an toàn phóng xạ, có bao nhiêu phát biểu đúng? Chỉ ghi một số nguyên. (1) Phóng xạ được dùng trong y học, công nghiệp và định tuổi; an toàn dựa trên giảm thời gian, tăng khoảng cách và che chắn phù hợp. (2) Cần đối chiếu kết quả với điều kiện bài toán. (3) Có thể tiếp xúc nguồn phóng xạ lâu nếu nguồn có kích thước nhỏ. (4) Mọi trường hợp đều cho cùng một kết quả.
\shortans{2}
\loigiai{
Đối chiếu từng phát biểu với kiến thức cốt lõi: Phóng xạ được dùng trong y học, công nghiệp và định tuổi; an toàn dựa trên giảm thời gian, tăng khoảng cách và che chắn phù hợp. Có 2 phát biểu đúng.
}
\end{ex}

% ===== Câu 132 =====
% ID: L12C4B23-34-TN
% Mức: TH
\dangbt{Tính chu kì bán rã, hằng số phóng xạ và tuổi mẫu vật}
\begin{ex}
% ID: L12C4B23-34-TN
% Mức: TH
Nội dung nào mô tả đúng nhất về tính chu kì bán rã, hằng số phóng xạ và tuổi mẫu vật?
\choice
{Chu kì bán rã càng lớn thì hằng số phóng xạ càng lớn.}
{Có thể bỏ qua điều kiện áp dụng và đơn vị trong mọi trường hợp.}
{Kết quả của bài toán không phụ thuộc dữ kiện đã cho.}
{\True Hằng số phóng xạ $\lambda=\ln2/T$; từ tỉ số lượng còn lại có thể suy ra thời gian hay tuổi mẫu.}
\loigiai{
Hằng số phóng xạ $\lambda=\ln2/T$; từ tỉ số lượng còn lại có thể suy ra thời gian hay tuổi mẫu. Vì vậy phương án $D$ đúng; các phương án còn lại trái với kiến thức hoặc điều kiện áp dụng.
}
\end{ex}

% ===== Câu 133 =====
% ID: L12C4B23-35-TL
% Mức: VD
\dangbt{Ứng dụng, ảnh hưởng và an toàn phóng xạ}
\begin{bt}
% ID: L12C4B23-35-TL
% Mức: VD
Nêu một tình huống thực tiễn liên quan đến ứng dụng, ảnh hưởng và an toàn phóng xạ và giải thích bằng kiến thức Vật lí.
\loigiai{
Cần nêu được: Phóng xạ được dùng trong y học, công nghiệp và định tuổi; an toàn dựa trên giảm thời gian, tăng khoảng cách và che chắn phù hợp. Khi vận dụng phải xác định đúng dữ kiện, điều kiện áp dụng, hệ đơn vị và kiểm tra tính hợp lí của kết quả. Nhận định “Có thể tiếp xúc nguồn phóng xạ lâu nếu nguồn có kích thước nhỏ.” sai vì trái với kiến thức trên.
}
\end{bt}

% ===== Câu 134 =====
% ID: L12C4B23-35-TLN
% Mức: VD
\begin{ex}
% ID: L12C4B23-35-TLN
% Mức: VD
Trong bốn phát biểu sau về ứng dụng, ảnh hưởng và an toàn phóng xạ, có bao nhiêu phát biểu đúng? Chỉ ghi một số nguyên. (1) Có thể tiếp xúc nguồn phóng xạ lâu nếu nguồn có kích thước nhỏ. (2) Cần đọc đúng dữ kiện và giả thiết. (3) Phóng xạ được dùng trong y học, công nghiệp và định tuổi; an toàn dựa trên giảm thời gian, tăng khoảng cách và che chắn phù hợp. (4) Có thể thay số trước khi chọn công thức.
\shortans{2}
\loigiai{
Đối chiếu từng phát biểu với kiến thức cốt lõi: Phóng xạ được dùng trong y học, công nghiệp và định tuổi; an toàn dựa trên giảm thời gian, tăng khoảng cách và che chắn phù hợp. Có 2 phát biểu đúng.
}
\end{ex}

% ===== Câu 135 =====
% ID: L12C4B23-35-TN
% Mức: TH
\dangbt{Tính chu kì bán rã, hằng số phóng xạ và tuổi mẫu vật}
\begin{ex}
% ID: L12C4B23-35-TN
% Mức: TH
Khi phân tích, phát biểu nào đúng về tính chu kì bán rã, hằng số phóng xạ và tuổi mẫu vật?
\choice
{\True Hằng số phóng xạ $\lambda=\ln2/T$; từ tỉ số lượng còn lại có thể suy ra thời gian hay tuổi mẫu.}
{Chu kì bán rã càng lớn thì hằng số phóng xạ càng lớn.}
{Có thể bỏ qua điều kiện áp dụng và đơn vị trong mọi trường hợp.}
{Kết quả của bài toán không phụ thuộc dữ kiện đã cho.}
\loigiai{
Hằng số phóng xạ $\lambda=\ln2/T$; từ tỉ số lượng còn lại có thể suy ra thời gian hay tuổi mẫu. Vì vậy phương án $A$ đúng; các phương án còn lại trái với kiến thức hoặc điều kiện áp dụng.
}
\end{ex}

% ===== Câu 136 =====
% ID: L12C4B23-36-TL
% Mức: VDC
\dangbt{Ứng dụng, ảnh hưởng và an toàn phóng xạ}
\begin{bt}
% ID: L12C4B23-36-TL
% Mức: VDC
So sánh nhận định đúng “Phóng xạ được dùng trong y học, công nghiệp và định tuổi; an toàn dựa trên giảm thời gian, tăng khoảng cách và che chắn phù hợp.” với nhận định sai “Có thể tiếp xúc nguồn phóng xạ lâu nếu nguồn có kích thước nhỏ.”; chỉ rõ căn cứ Vật lí.
\loigiai{
Cần nêu được: Phóng xạ được dùng trong y học, công nghiệp và định tuổi; an toàn dựa trên giảm thời gian, tăng khoảng cách và che chắn phù hợp. Khi vận dụng phải xác định đúng dữ kiện, điều kiện áp dụng, hệ đơn vị và kiểm tra tính hợp lí của kết quả. Nhận định “Có thể tiếp xúc nguồn phóng xạ lâu nếu nguồn có kích thước nhỏ.” sai vì trái với kiến thức trên.
}
\end{bt}

% ===== Câu 137 =====
% ID: L12C4B23-36-TLN
% Mức: VDC
\begin{ex}
% ID: L12C4B23-36-TLN
% Mức: VDC
Trong bốn phát biểu sau về ứng dụng, ảnh hưởng và an toàn phóng xạ, có bao nhiêu phát biểu đúng? Chỉ ghi một số nguyên. (1) Kết quả cần có ý nghĩa vật lí. (2) Có thể tiếp xúc nguồn phóng xạ lâu nếu nguồn có kích thước nhỏ. (3) Phải dùng đúng định luật cho tình huống. (4) Phóng xạ được dùng trong y học, công nghiệp và định tuổi; an toàn dựa trên giảm thời gian, tăng khoảng cách và che chắn phù hợp.
\shortans{3}
\loigiai{
Đối chiếu từng phát biểu với kiến thức cốt lõi: Phóng xạ được dùng trong y học, công nghiệp và định tuổi; an toàn dựa trên giảm thời gian, tăng khoảng cách và che chắn phù hợp. Có 3 phát biểu đúng.
}
\end{ex}

% ===== Câu 138 =====
% ID: L12C4B23-36-TN
% Mức: TH
\dangbt{Tính chu kì bán rã, hằng số phóng xạ và tuổi mẫu vật}
\begin{ex}
% ID: L12C4B23-36-TN
% Mức: TH
Kết luận nào của học sinh là đúng về tính chu kì bán rã, hằng số phóng xạ và tuổi mẫu vật?
\choice
{Chu kì bán rã càng lớn thì hằng số phóng xạ càng lớn.}
{\True Hằng số phóng xạ $\lambda=\ln2/T$; từ tỉ số lượng còn lại có thể suy ra thời gian hay tuổi mẫu.}
{Có thể bỏ qua điều kiện áp dụng và đơn vị trong mọi trường hợp.}
{Kết quả của bài toán không phụ thuộc dữ kiện đã cho.}
\loigiai{
Hằng số phóng xạ $\lambda=\ln2/T$; từ tỉ số lượng còn lại có thể suy ra thời gian hay tuổi mẫu. Vì vậy phương án $B$ đúng; các phương án còn lại trái với kiến thức hoặc điều kiện áp dụng.
}
\end{ex}

% ===== Câu 139 =====
% ID: L12C4B23-37-TN
% Mức: TH
\begin{ex}
% ID: L12C4B23-37-TN
% Mức: TH
Chọn phát biểu không trái với kiến thức về tính chu kì bán rã, hằng số phóng xạ và tuổi mẫu vật?
\choice
{Chu kì bán rã càng lớn thì hằng số phóng xạ càng lớn.}
{Có thể bỏ qua điều kiện áp dụng và đơn vị trong mọi trường hợp.}
{\True Hằng số phóng xạ $\lambda=\ln2/T$; từ tỉ số lượng còn lại có thể suy ra thời gian hay tuổi mẫu.}
{Kết quả của bài toán không phụ thuộc dữ kiện đã cho.}
\loigiai{
Hằng số phóng xạ $\lambda=\ln2/T$; từ tỉ số lượng còn lại có thể suy ra thời gian hay tuổi mẫu. Vì vậy phương án $C$ đúng; các phương án còn lại trái với kiến thức hoặc điều kiện áp dụng.
}
\end{ex}

% ===== Câu 140 =====
% ID: L12C4B23-38-TN
% Mức: VD
\begin{ex}
% ID: L12C4B23-38-TN
% Mức: VD
Cơ sở Vật lí đúng để giải bài là gì về tính chu kì bán rã, hằng số phóng xạ và tuổi mẫu vật?
\choice
{Chu kì bán rã càng lớn thì hằng số phóng xạ càng lớn.}
{Có thể bỏ qua điều kiện áp dụng và đơn vị trong mọi trường hợp.}
{Kết quả của bài toán không phụ thuộc dữ kiện đã cho.}
{\True Hằng số phóng xạ $\lambda=\ln2/T$; từ tỉ số lượng còn lại có thể suy ra thời gian hay tuổi mẫu.}
\loigiai{
Hằng số phóng xạ $\lambda=\ln2/T$; từ tỉ số lượng còn lại có thể suy ra thời gian hay tuổi mẫu. Vì vậy phương án $D$ đúng; các phương án còn lại trái với kiến thức hoặc điều kiện áp dụng.
}
\end{ex}

% ===== Câu 141 =====
% ID: L12C4B23-39-TN
% Mức: VD
\begin{ex}
% ID: L12C4B23-39-TN
% Mức: VD
Trong một bài thực hành, nhận xét nào đúng về tính chu kì bán rã, hằng số phóng xạ và tuổi mẫu vật?
\choice
{\True Hằng số phóng xạ $\lambda=\ln2/T$; từ tỉ số lượng còn lại có thể suy ra thời gian hay tuổi mẫu.}
{Chu kì bán rã càng lớn thì hằng số phóng xạ càng lớn.}
{Có thể bỏ qua điều kiện áp dụng và đơn vị trong mọi trường hợp.}
{Kết quả của bài toán không phụ thuộc dữ kiện đã cho.}
\loigiai{
Hằng số phóng xạ $\lambda=\ln2/T$; từ tỉ số lượng còn lại có thể suy ra thời gian hay tuổi mẫu. Vì vậy phương án $A$ đúng; các phương án còn lại trái với kiến thức hoặc điều kiện áp dụng.
}
\end{ex}

% ===== Câu 142 =====
% ID: L12C4B23-40-TN
% Mức: VD
\begin{ex}
% ID: L12C4B23-40-TN
% Mức: VD
Kết luận đúng khi vận dụng là về tính chu kì bán rã, hằng số phóng xạ và tuổi mẫu vật?
\choice
{Chu kì bán rã càng lớn thì hằng số phóng xạ càng lớn.}
{\True Hằng số phóng xạ $\lambda=\ln2/T$; từ tỉ số lượng còn lại có thể suy ra thời gian hay tuổi mẫu.}
{Có thể bỏ qua điều kiện áp dụng và đơn vị trong mọi trường hợp.}
{Kết quả của bài toán không phụ thuộc dữ kiện đã cho.}
\loigiai{
Hằng số phóng xạ $\lambda=\ln2/T$; từ tỉ số lượng còn lại có thể suy ra thời gian hay tuổi mẫu. Vì vậy phương án $B$ đúng; các phương án còn lại trái với kiến thức hoặc điều kiện áp dụng.
}
\end{ex}

% ===== Câu 143 =====
% ID: L12C4B23-41-TN
% Mức: NB
\dangbt{Tính độ phóng xạ và khai thác đồ thị phóng xạ}
\begin{ex}
% ID: L12C4B23-41-TN
% Mức: NB
Phát biểu nào sau đây đúng về tính độ phóng xạ và khai thác đồ thị phóng xạ?
\choice
{\True Độ phóng xạ $H=\lambda N$ và cũng giảm theo quy luật mũ với cùng chu kì bán rã.}
{Độ phóng xạ không phụ thuộc số hạt chưa phân rã.}
{Có thể bỏ qua điều kiện áp dụng và đơn vị trong mọi trường hợp.}
{Kết quả của bài toán không phụ thuộc dữ kiện đã cho.}
\loigiai{
Độ phóng xạ $H=\lambda N$ và cũng giảm theo quy luật mũ với cùng chu kì bán rã. Vì vậy phương án $A$ đúng; các phương án còn lại trái với kiến thức hoặc điều kiện áp dụng.
}
\end{ex}

% ===== Câu 144 =====
% ID: L12C4B23-42-TN
% Mức: NB
\begin{ex}
% ID: L12C4B23-42-TN
% Mức: NB
Nhận định nào phù hợp với kiến thức Vật lí về tính độ phóng xạ và khai thác đồ thị phóng xạ?
\choice
{Độ phóng xạ không phụ thuộc số hạt chưa phân rã.}
{\True Độ phóng xạ $H=\lambda N$ và cũng giảm theo quy luật mũ với cùng chu kì bán rã.}
{Có thể bỏ qua điều kiện áp dụng và đơn vị trong mọi trường hợp.}
{Kết quả của bài toán không phụ thuộc dữ kiện đã cho.}
\loigiai{
Độ phóng xạ $H=\lambda N$ và cũng giảm theo quy luật mũ với cùng chu kì bán rã. Vì vậy phương án $B$ đúng; các phương án còn lại trái với kiến thức hoặc điều kiện áp dụng.
}
\end{ex}

% ===== Câu 145 =====
% ID: L12C4B23-43-TN
% Mức: NB
\begin{ex}
% ID: L12C4B23-43-TN
% Mức: NB
Chọn kết luận chính xác về tính độ phóng xạ và khai thác đồ thị phóng xạ?
\choice
{Độ phóng xạ không phụ thuộc số hạt chưa phân rã.}
{Có thể bỏ qua điều kiện áp dụng và đơn vị trong mọi trường hợp.}
{\True Độ phóng xạ $H=\lambda N$ và cũng giảm theo quy luật mũ với cùng chu kì bán rã.}
{Kết quả của bài toán không phụ thuộc dữ kiện đã cho.}
\loigiai{
Độ phóng xạ $H=\lambda N$ và cũng giảm theo quy luật mũ với cùng chu kì bán rã. Vì vậy phương án $C$ đúng; các phương án còn lại trái với kiến thức hoặc điều kiện áp dụng.
}
\end{ex}

% ===== Câu 146 =====
% ID: L12C4B23-44-TN
% Mức: TH
\begin{ex}
% ID: L12C4B23-44-TN
% Mức: TH
Nội dung nào mô tả đúng nhất về tính độ phóng xạ và khai thác đồ thị phóng xạ?
\choice
{Độ phóng xạ không phụ thuộc số hạt chưa phân rã.}
{Có thể bỏ qua điều kiện áp dụng và đơn vị trong mọi trường hợp.}
{Kết quả của bài toán không phụ thuộc dữ kiện đã cho.}
{\True Độ phóng xạ $H=\lambda N$ và cũng giảm theo quy luật mũ với cùng chu kì bán rã.}
\loigiai{
Độ phóng xạ $H=\lambda N$ và cũng giảm theo quy luật mũ với cùng chu kì bán rã. Vì vậy phương án $D$ đúng; các phương án còn lại trái với kiến thức hoặc điều kiện áp dụng.
}
\end{ex}

% ===== Câu 147 =====
% ID: L12C4B23-45-TN
% Mức: TH
\begin{ex}
% ID: L12C4B23-45-TN
% Mức: TH
Khi phân tích, phát biểu nào đúng về tính độ phóng xạ và khai thác đồ thị phóng xạ?
\choice
{\True Độ phóng xạ $H=\lambda N$ và cũng giảm theo quy luật mũ với cùng chu kì bán rã.}
{Độ phóng xạ không phụ thuộc số hạt chưa phân rã.}
{Có thể bỏ qua điều kiện áp dụng và đơn vị trong mọi trường hợp.}
{Kết quả của bài toán không phụ thuộc dữ kiện đã cho.}
\loigiai{
Độ phóng xạ $H=\lambda N$ và cũng giảm theo quy luật mũ với cùng chu kì bán rã. Vì vậy phương án $A$ đúng; các phương án còn lại trái với kiến thức hoặc điều kiện áp dụng.
}
\end{ex}

% ===== Câu 148 =====
% ID: L12C4B23-46-TN
% Mức: TH
\begin{ex}
% ID: L12C4B23-46-TN
% Mức: TH
Kết luận nào của học sinh là đúng về tính độ phóng xạ và khai thác đồ thị phóng xạ?
\choice
{Độ phóng xạ không phụ thuộc số hạt chưa phân rã.}
{\True Độ phóng xạ $H=\lambda N$ và cũng giảm theo quy luật mũ với cùng chu kì bán rã.}
{Có thể bỏ qua điều kiện áp dụng và đơn vị trong mọi trường hợp.}
{Kết quả của bài toán không phụ thuộc dữ kiện đã cho.}
\loigiai{
Độ phóng xạ $H=\lambda N$ và cũng giảm theo quy luật mũ với cùng chu kì bán rã. Vì vậy phương án $B$ đúng; các phương án còn lại trái với kiến thức hoặc điều kiện áp dụng.
}
\end{ex}

% ===== Câu 149 =====
% ID: L12C4B23-47-TN
% Mức: TH
\begin{ex}
% ID: L12C4B23-47-TN
% Mức: TH
Chọn phát biểu không trái với kiến thức về tính độ phóng xạ và khai thác đồ thị phóng xạ?
\choice
{Độ phóng xạ không phụ thuộc số hạt chưa phân rã.}
{Có thể bỏ qua điều kiện áp dụng và đơn vị trong mọi trường hợp.}
{\True Độ phóng xạ $H=\lambda N$ và cũng giảm theo quy luật mũ với cùng chu kì bán rã.}
{Kết quả của bài toán không phụ thuộc dữ kiện đã cho.}
\loigiai{
Độ phóng xạ $H=\lambda N$ và cũng giảm theo quy luật mũ với cùng chu kì bán rã. Vì vậy phương án $C$ đúng; các phương án còn lại trái với kiến thức hoặc điều kiện áp dụng.
}
\end{ex}

% ===== Câu 150 =====
% ID: L12C4B23-48-TN
% Mức: VD
\begin{ex}
% ID: L12C4B23-48-TN
% Mức: VD
Cơ sở Vật lí đúng để giải bài là gì về tính độ phóng xạ và khai thác đồ thị phóng xạ?
\choice
{Độ phóng xạ không phụ thuộc số hạt chưa phân rã.}
{Có thể bỏ qua điều kiện áp dụng và đơn vị trong mọi trường hợp.}
{Kết quả của bài toán không phụ thuộc dữ kiện đã cho.}
{\True Độ phóng xạ $H=\lambda N$ và cũng giảm theo quy luật mũ với cùng chu kì bán rã.}
\loigiai{
Độ phóng xạ $H=\lambda N$ và cũng giảm theo quy luật mũ với cùng chu kì bán rã. Vì vậy phương án $D$ đúng; các phương án còn lại trái với kiến thức hoặc điều kiện áp dụng.
}
\end{ex}

% ===== Câu 151 =====
% ID: L12C4B23-49-TN
% Mức: VD
\begin{ex}
% ID: L12C4B23-49-TN
% Mức: VD
Trong một bài thực hành, nhận xét nào đúng về tính độ phóng xạ và khai thác đồ thị phóng xạ?
\choice
{\True Độ phóng xạ $H=\lambda N$ và cũng giảm theo quy luật mũ với cùng chu kì bán rã.}
{Độ phóng xạ không phụ thuộc số hạt chưa phân rã.}
{Có thể bỏ qua điều kiện áp dụng và đơn vị trong mọi trường hợp.}
{Kết quả của bài toán không phụ thuộc dữ kiện đã cho.}
\loigiai{
Độ phóng xạ $H=\lambda N$ và cũng giảm theo quy luật mũ với cùng chu kì bán rã. Vì vậy phương án $A$ đúng; các phương án còn lại trái với kiến thức hoặc điều kiện áp dụng.
}
\end{ex}

% ===== Câu 152 =====
% ID: L12C4B23-50-TN
% Mức: VD
\begin{ex}
% ID: L12C4B23-50-TN
% Mức: VD
Kết luận đúng khi vận dụng là về tính độ phóng xạ và khai thác đồ thị phóng xạ?
\choice
{Độ phóng xạ không phụ thuộc số hạt chưa phân rã.}
{\True Độ phóng xạ $H=\lambda N$ và cũng giảm theo quy luật mũ với cùng chu kì bán rã.}
{Có thể bỏ qua điều kiện áp dụng và đơn vị trong mọi trường hợp.}
{Kết quả của bài toán không phụ thuộc dữ kiện đã cho.}
\loigiai{
Độ phóng xạ $H=\lambda N$ và cũng giảm theo quy luật mũ với cùng chu kì bán rã. Vì vậy phương án $B$ đúng; các phương án còn lại trái với kiến thức hoặc điều kiện áp dụng.
}
\end{ex}

% ===== Câu 153 =====
% ID: L12C4B23-51-TN
% Mức: NB
\dangbt{Ứng dụng, ảnh hưởng và an toàn phóng xạ}
\begin{ex}
% ID: L12C4B23-51-TN
% Mức: NB
Phát biểu nào sau đây đúng về ứng dụng, ảnh hưởng và an toàn phóng xạ?
\choice
{\True Phóng xạ được dùng trong y học, công nghiệp và định tuổi; an toàn dựa trên giảm thời gian, tăng khoảng cách và che chắn phù hợp.}
{Có thể tiếp xúc nguồn phóng xạ lâu nếu nguồn có kích thước nhỏ.}
{Có thể bỏ qua điều kiện áp dụng và đơn vị trong mọi trường hợp.}
{Kết quả của bài toán không phụ thuộc dữ kiện đã cho.}
\loigiai{
Phóng xạ được dùng trong y học, công nghiệp và định tuổi; an toàn dựa trên giảm thời gian, tăng khoảng cách và che chắn phù hợp. Vì vậy phương án $A$ đúng; các phương án còn lại trái với kiến thức hoặc điều kiện áp dụng.
}
\end{ex}

% ===== Câu 154 =====
% ID: L12C4B23-52-TN
% Mức: NB
\begin{ex}
% ID: L12C4B23-52-TN
% Mức: NB
Nhận định nào phù hợp với kiến thức Vật lí về ứng dụng, ảnh hưởng và an toàn phóng xạ?
\choice
{Có thể tiếp xúc nguồn phóng xạ lâu nếu nguồn có kích thước nhỏ.}
{\True Phóng xạ được dùng trong y học, công nghiệp và định tuổi; an toàn dựa trên giảm thời gian, tăng khoảng cách và che chắn phù hợp.}
{Có thể bỏ qua điều kiện áp dụng và đơn vị trong mọi trường hợp.}
{Kết quả của bài toán không phụ thuộc dữ kiện đã cho.}
\loigiai{
Phóng xạ được dùng trong y học, công nghiệp và định tuổi; an toàn dựa trên giảm thời gian, tăng khoảng cách và che chắn phù hợp. Vì vậy phương án $B$ đúng; các phương án còn lại trái với kiến thức hoặc điều kiện áp dụng.
}
\end{ex}

% ===== Câu 155 =====
% ID: L12C4B23-53-TN
% Mức: NB
\begin{ex}
% ID: L12C4B23-53-TN
% Mức: NB
Chọn kết luận chính xác về ứng dụng, ảnh hưởng và an toàn phóng xạ?
\choice
{Có thể tiếp xúc nguồn phóng xạ lâu nếu nguồn có kích thước nhỏ.}
{Có thể bỏ qua điều kiện áp dụng và đơn vị trong mọi trường hợp.}
{\True Phóng xạ được dùng trong y học, công nghiệp và định tuổi; an toàn dựa trên giảm thời gian, tăng khoảng cách và che chắn phù hợp.}
{Kết quả của bài toán không phụ thuộc dữ kiện đã cho.}
\loigiai{
Phóng xạ được dùng trong y học, công nghiệp và định tuổi; an toàn dựa trên giảm thời gian, tăng khoảng cách và che chắn phù hợp. Vì vậy phương án $C$ đúng; các phương án còn lại trái với kiến thức hoặc điều kiện áp dụng.
}
\end{ex}

% ===== Câu 156 =====
% ID: L12C4B23-54-TN
% Mức: TH
\begin{ex}
% ID: L12C4B23-54-TN
% Mức: TH
Nội dung nào mô tả đúng nhất về ứng dụng, ảnh hưởng và an toàn phóng xạ?
\choice
{Có thể tiếp xúc nguồn phóng xạ lâu nếu nguồn có kích thước nhỏ.}
{Có thể bỏ qua điều kiện áp dụng và đơn vị trong mọi trường hợp.}
{Kết quả của bài toán không phụ thuộc dữ kiện đã cho.}
{\True Phóng xạ được dùng trong y học, công nghiệp và định tuổi; an toàn dựa trên giảm thời gian, tăng khoảng cách và che chắn phù hợp.}
\loigiai{
Phóng xạ được dùng trong y học, công nghiệp và định tuổi; an toàn dựa trên giảm thời gian, tăng khoảng cách và che chắn phù hợp. Vì vậy phương án $D$ đúng; các phương án còn lại trái với kiến thức hoặc điều kiện áp dụng.
}
\end{ex}

% ===== Câu 157 =====
% ID: L12C4B23-55-TN
% Mức: TH
\begin{ex}
% ID: L12C4B23-55-TN
% Mức: TH
Khi phân tích, phát biểu nào đúng về ứng dụng, ảnh hưởng và an toàn phóng xạ?
\choice
{\True Phóng xạ được dùng trong y học, công nghiệp và định tuổi; an toàn dựa trên giảm thời gian, tăng khoảng cách và che chắn phù hợp.}
{Có thể tiếp xúc nguồn phóng xạ lâu nếu nguồn có kích thước nhỏ.}
{Có thể bỏ qua điều kiện áp dụng và đơn vị trong mọi trường hợp.}
{Kết quả của bài toán không phụ thuộc dữ kiện đã cho.}
\loigiai{
Phóng xạ được dùng trong y học, công nghiệp và định tuổi; an toàn dựa trên giảm thời gian, tăng khoảng cách và che chắn phù hợp. Vì vậy phương án $A$ đúng; các phương án còn lại trái với kiến thức hoặc điều kiện áp dụng.
}
\end{ex}

% ===== Câu 158 =====
% ID: L12C4B23-56-TN
% Mức: TH
\begin{ex}
% ID: L12C4B23-56-TN
% Mức: TH
Kết luận nào của học sinh là đúng về ứng dụng, ảnh hưởng và an toàn phóng xạ?
\choice
{Có thể tiếp xúc nguồn phóng xạ lâu nếu nguồn có kích thước nhỏ.}
{\True Phóng xạ được dùng trong y học, công nghiệp và định tuổi; an toàn dựa trên giảm thời gian, tăng khoảng cách và che chắn phù hợp.}
{Có thể bỏ qua điều kiện áp dụng và đơn vị trong mọi trường hợp.}
{Kết quả của bài toán không phụ thuộc dữ kiện đã cho.}
\loigiai{
Phóng xạ được dùng trong y học, công nghiệp và định tuổi; an toàn dựa trên giảm thời gian, tăng khoảng cách và che chắn phù hợp. Vì vậy phương án $B$ đúng; các phương án còn lại trái với kiến thức hoặc điều kiện áp dụng.
}
\end{ex}

% ===== Câu 159 =====
% ID: L12C4B23-57-TN
% Mức: TH
\begin{ex}
% ID: L12C4B23-57-TN
% Mức: TH
Chọn phát biểu không trái với kiến thức về ứng dụng, ảnh hưởng và an toàn phóng xạ?
\choice
{Có thể tiếp xúc nguồn phóng xạ lâu nếu nguồn có kích thước nhỏ.}
{Có thể bỏ qua điều kiện áp dụng và đơn vị trong mọi trường hợp.}
{\True Phóng xạ được dùng trong y học, công nghiệp và định tuổi; an toàn dựa trên giảm thời gian, tăng khoảng cách và che chắn phù hợp.}
{Kết quả của bài toán không phụ thuộc dữ kiện đã cho.}
\loigiai{
Phóng xạ được dùng trong y học, công nghiệp và định tuổi; an toàn dựa trên giảm thời gian, tăng khoảng cách và che chắn phù hợp. Vì vậy phương án $C$ đúng; các phương án còn lại trái với kiến thức hoặc điều kiện áp dụng.
}
\end{ex}

% ===== Câu 160 =====
% ID: L12C4B23-58-TN
% Mức: VD
\begin{ex}
% ID: L12C4B23-58-TN
% Mức: VD
Cơ sở Vật lí đúng để giải bài là gì về ứng dụng, ảnh hưởng và an toàn phóng xạ?
\choice
{Có thể tiếp xúc nguồn phóng xạ lâu nếu nguồn có kích thước nhỏ.}
{Có thể bỏ qua điều kiện áp dụng và đơn vị trong mọi trường hợp.}
{Kết quả của bài toán không phụ thuộc dữ kiện đã cho.}
{\True Phóng xạ được dùng trong y học, công nghiệp và định tuổi; an toàn dựa trên giảm thời gian, tăng khoảng cách và che chắn phù hợp.}
\loigiai{
Phóng xạ được dùng trong y học, công nghiệp và định tuổi; an toàn dựa trên giảm thời gian, tăng khoảng cách và che chắn phù hợp. Vì vậy phương án $D$ đúng; các phương án còn lại trái với kiến thức hoặc điều kiện áp dụng.
}
\end{ex}

% ===== Câu 161 =====
% ID: L12C4B23-59-TN
% Mức: VD
\begin{ex}
% ID: L12C4B23-59-TN
% Mức: VD
Trong một bài thực hành, nhận xét nào đúng về ứng dụng, ảnh hưởng và an toàn phóng xạ?
\choice
{\True Phóng xạ được dùng trong y học, công nghiệp và định tuổi; an toàn dựa trên giảm thời gian, tăng khoảng cách và che chắn phù hợp.}
{Có thể tiếp xúc nguồn phóng xạ lâu nếu nguồn có kích thước nhỏ.}
{Có thể bỏ qua điều kiện áp dụng và đơn vị trong mọi trường hợp.}
{Kết quả của bài toán không phụ thuộc dữ kiện đã cho.}
\loigiai{
Phóng xạ được dùng trong y học, công nghiệp và định tuổi; an toàn dựa trên giảm thời gian, tăng khoảng cách và che chắn phù hợp. Vì vậy phương án $A$ đúng; các phương án còn lại trái với kiến thức hoặc điều kiện áp dụng.
}
\end{ex}

% ===== Câu 162 =====
% ID: L12C4B23-60-TN
% Mức: VD
\begin{ex}
% ID: L12C4B23-60-TN
% Mức: VD
Kết luận đúng khi vận dụng là về ứng dụng, ảnh hưởng và an toàn phóng xạ?
\choice
{Có thể tiếp xúc nguồn phóng xạ lâu nếu nguồn có kích thước nhỏ.}
{\True Phóng xạ được dùng trong y học, công nghiệp và định tuổi; an toàn dựa trên giảm thời gian, tăng khoảng cách và che chắn phù hợp.}
{Có thể bỏ qua điều kiện áp dụng và đơn vị trong mọi trường hợp.}
{Kết quả của bài toán không phụ thuộc dữ kiện đã cho.}
\loigiai{
Phóng xạ được dùng trong y học, công nghiệp và định tuổi; an toàn dựa trên giảm thời gian, tăng khoảng cách và che chắn phù hợp. Vì vậy phương án $B$ đúng; các phương án còn lại trái với kiến thức hoặc điều kiện áp dụng.
}
\end{ex}
